% Section 3, part 1 of 7; label prefix: gau. The main file inputs only this part; each part inputs the next.
% Part 1: introduction; the cyclic form of PF4 (objects, constants, self-adjointness, the concomitant).
% Part 2: transversality, isotropy, the conifold expansion, special geometry, (CYC), the cyclic system.
% Part 3: the quintic system and the pulled-back model; Theorem U (statement), the commutant.
% Part 4: lifting, the scalar equation, proof of Theorem U, the leading term, the Z/5-degree, graded gauges.
% Part 5: (R1), the diagonal, the vanishing solutions, the gauge identity (G), the canonical gauge, symplecticity.
% Part 6: HPF, polynomiality, formal WKB uniqueness, colour covariance of the model columns.
% Part 7: colour covariance of the actual rows, the normalisation series D, the rows in the ring, the colours of D, (F_g).
% No new macros. Local letters are declared where they occur.
\section{Cyclic systems, Theorem U and the canonical gauge}
\label{sec:gau}

This section relates the cyclic system of the formal quintic at the conifold point to that of the model and defines the normalisation series $\Dn$ by which their formal solutions differ. \Cref{sec:gau-cyc} proves the cyclic form (CYC) of $\PFfour$ in every conifold frame, so that the cyclic system closes on $\PFfive$. \Cref{sec:gau-systems} writes the quintic system and the model pulled back along $v(u)$ as cyclic systems with one simple pole and equal products. \Cref{sec:gau-U} proves Theorem~U over every field of characteristic $0$: the gauges between two such systems form a free module of rank one over the constant series in $z$. \Cref{sec:gau-van} proves (R1), computes the diagonal of a gauge and proves the gauge identity (G) between the vanishing solutions, which singles out the canonical gauge $\Pc$; \cref{sec:gau-symp} proves that $\Pc$ is symplectic. \Cref{sec:gau-R0} compares $\Pc$ applied to the model columns with the rows of \cref{def:set-rows}: they differ by a constant diagonal series $\Dn=\diag(\Dn_i)$ with $\Dn_i(z)=\Dn_0(\zeta^{-i}z)$.

\emph{Conventions.} $\Dq=q\,\dd/\dd q=-(1-u)\,\dd/\dd u$ (\cref{def:set-genus0,def:set-points}); for a series in $u$ a prime denotes $\dd/\dd u$; $z$ is a constant, and $z\de_u$ is $z\,\dd/\dd u$ at fixed $z$, so $z\Dq=-(1-u)z\de_u$. $E_{ab}$ ($a,b\in\bbZ/5$) are the matrix units and $\mathbf e_a$ the standard basis vectors; solutions, gauges and rows are in lowered components (\cref{def:set-rows}). The quoted inputs enter this section only through \cref{sec:set}: (M) and (G0) through \cref{prop:set-cyclicLR,cor:set-cyclicPF5}, and (QRR) only in Step~3 of the proof of \cref{lem:gau-colrows}. Everything else is proved here, by identities of formal series. Finite computations testing this section are reported in \cref{sec:evi-gauge}; no proof uses them.

%-------------------------------------------------------------------------------
\subsection{The cyclic form of the fourth-order operator}
\label{sec:gau-cyc}

In this subsection $\PFfour=\sum_{k=0}^4A_k\Dq^k$ is written in the variable $u$ (\cref{def:set-genus0}):
\[
A_4=u,\qquad A_3=-2(1-u),\qquad A_2=-\tfrac75(1-u),\qquad A_1=-\tfrac25(1-u),\qquad A_0=-\tfrac{24}{625}(1-u),
\]
and we use $\Dq(1-u)=1-u$ and $\Dq u=-(1-u)$ throughout.

\begin{defn}[Objects of the cyclic-form theorem]
\label{def:gau-conc}
\begin{enumerate}
\item $\bbK:=\bbQ\lser{u}$ with the derivation $\Dq$, and $\bbK[\ell]$ is the polynomial ring in an indeterminate $\ell$ with the unique extension of $\Dq$ for which $\Dq\ell=1-u^{-1}$ (so $\ell$ plays the role of $\log u$). $\bbK[\Dq]$ is the ring of differential operators $\sum_ka_k\Dq^k$, $a_k\in\bbK$, with $\Dq\circ a-a\circ\Dq=\Dq(a)$; it acts on every differential ring that extends $(\bbK,\Dq)$.
\item $f_0,f_1,f_2\in\bbQ\fps u$ are the solutions of $\PFfour$ with $f_k=u^k+O(u^3)$, and $\varpi_1=f_1+\gamma f_2$, $\gamma=\frac7{10}$, is the vanishing period (\cref{lem:set-frobbasis}). $y_3:=\varpi_1\ell+\rho_3\in\bbK[\ell]$, $\rho_3\in u^3\bbQ\fps u$, is the logarithmic solution of \cref{lem:set-frobbasis}~(iii), whose indeterminate $\ell$ is the one above; thus $\PFfour y_3=0$ in $\bbK[\ell]$. Put $\mathcal S:=\operatorname{span}_\bbQ(f_0,f_1,f_2,y_3)\subset\bbK[\ell]$; no claim is made that $\mathcal S$ contains every solution in $\bbK[\ell]$. An element of $\mathcal S$ lies in $\bbK$ if and only if its coefficient of $y_3$ vanishes, so $\mathcal S\cap\bbK=\operatorname{span}(f_0,f_1,f_2)$, the space of power-series solutions of $\PFfour$.
\item For $\cP=\sum_{k\leq4}a_k\Dq^k$ with coefficients $a_k$ in a differential ring $(\cR_0,\Dq)$, for instance $\cP\in\bbK[\Dq]$, the \emph{formal adjoint} is $\cP^*:=\sum_k(-\Dq)^k\circ a_k$, and the \emph{concomitant} of $\cP$ is the bilinear differential expression
\[
\Conc_\cP(f,g):=\sum_{k=1}^4\sum_{j=0}^{k-1}(-1)^j\,\Dq^j(a_kg)\,\Dq^{k-1-j}f .
\]
We write $\Conc:=\Conc_{\PFfour}$ and expand it as $\Conc(f,g)=\sum_{a,b=0}^3\Dq^af\,\Jet_{ab}\,\Dq^bg$ with $\Jet_{ab}\in\bbQ[u]\subset\bbK$; $\Jet=(\Jet_{ab})$ is the \emph{jet matrix}. The $A_k$ are polynomials in $u$, so $\PFfour$, $\PFfour^*$ and $\Conc$ have coefficients in the differential ring $\bbQ[u]$, $\Dq u=-(1-u)$, and act on every differential ring containing it: on $\bbK[\ell]$, and also on $\bbQ\fps q[\logq]$ (with $u=1-5^5q$), on $\Nring$ (with $u=L^{-5}$) and on $\bbQ\lser s$, where they are used in \cref{def:set-confexp,lem:gau-confexp,cor:gau-rowsNN}.
\item For a basis $\mathbf x=(x_0,\dots,x_3)$ of $\mathcal S$, the \emph{Gram matrix} is $\Conc_{\mathbf x}:=(\Conc(x_i,x_j))_{i,j=0}^3$ and the \emph{Wronskian matrix} is $\Wr_{\mathbf x}:=(\Dq^ax_i)_{a,i=0}^3$ (row $a$, column $i$), whose determinant is the Wronskian $\Wr(x_0,\dots,x_3)$ in the sense of \cref{lem:set-wronskian}. The basis $\mathbf f:=(f_0,f_1,f_2,y_3)$ is the standard one.
\item $\sh$ is the $\bbK$-algebra automorphism of $\bbK[\ell]$ with $\sh(\ell)=\ell+1$.
\item For $t\in\bbK$ with $\Dq t\neq0$, $\de_t:=(\Dq t)^{-1}\Dq$, a derivation of $\bbK[\ell]$. For a basis $\mathbf x$ of $\mathcal S$ (with $\Conc_{\mathbf x}$ invertible, \cref{lem:gau-conc}) the \emph{transversality sums} are $\Upsilon_k:=\sum_{i,j}(\Conc_{\mathbf x}^{-1})_{ij}\,x_i\,\Dq^kx_j$ and, for $y_0\in\bbK^\times$, $\widetilde\Upsilon_k:=\sum_{i,j}(\Conc_{\mathbf x}^{-1})_{ij}\,(x_i/y_0)\,\de_t^k(x_j/y_0)$ ($k\geq0$), elements of the fraction field of $\bbK[\ell]$.
\item For $y_0,y_1\in\bbK$ with $y_0\neq0$ and $t:=y_1/y_0$ not constant, the \emph{cyclic operator} is
\begin{equation}
\label{eq:gau-Mcyc}
\Mcyc(y_0,y_1):=\frac1{y_0}\circ\Dq\circ\frac1{\Dq t}\circ\Dq\circ\bigl(u\,y_0^2(\Dq t)^2\bigr)\circ\Dq\circ\frac1{\Dq t}\circ\Dq\circ\frac1{y_0}\ \in\bbK[\Dq].
\end{equation}
\end{enumerate}
\end{defn}

\begin{lem}[Constants of $\bbK{[\ell]}$]
\label{lem:gau-const}
$\{f\in\bbK[\ell]:\Dq f=0\}=\bbQ$.
\end{lem}

\begin{proof}
For $f\in\bbK$, $\Dq f=-(1-u)f'=0$ gives $f'=0$, since $1-u$ is a unit of $\bbK$; so $f\in\bbQ$, the characteristic being $0$. Let $f=\sum_{k\leq n}a_k\ell^k$ with $a_k\in\bbK$, $a_n\neq0$ and $n\geq1$, and suppose $\Dq f=0$. Since $\Dq f=\sum_k\bigl(\Dq a_k+(k+1)a_{k+1}\Dq\ell\bigr)\ell^k$, the coefficient of $\ell^n$ gives $\Dq a_n=0$, so $a_n\in\bbQ^\times$, and the coefficient of $\ell^{n-1}$ gives $\Dq a_{n-1}=-na_n\Dq\ell=na_n(1-u)/u$. Dividing by the unit $-(1-u)$ gives $a_{n-1}'=-na_n/u$. This is impossible: the derivative of a Laurent series has no term in $u^{-1}$, and $na_n\neq0$. Hence $n=0$.
\end{proof}

\begin{lem}[$\PFfour$ is self-adjoint]
\label{lem:gau-adj}
$\PFfour^*=\PFfour$, as operators with coefficients in $\bbQ[u]$; hence as operators on every differential ring containing $\bbQ[u]$.
\end{lem}

\begin{proof}
The formal adjoint is the anti-automorphism of $\bbQ[u][\Dq]$ (and of $\bbK[\Dq]$) that fixes the coefficients and sends $\Dq$ to $-\Dq$; it exists because it maps the relation $\Dq\circ a-a\circ\Dq=\Dq(a)$ to $a\circ(-\Dq)-(-\Dq)\circ a=\Dq(a)$, which holds. Here $q=(1-u)/5^5\in\bbQ[u]$ and $\Dq q=q$. Write $\PFfour=\Dq^4-5^5q\circ P_4(\Dq)$ with $P_4(\Dq):=\prod_{k=1}^4(\Dq+k/5)$, since $5q\prod_{k=1}^4(5\Dq+k)=5^5q\prod_{k=1}^4(\Dq+k/5)$. Then $(\Dq^4)^*=\Dq^4$ and $(q\circ P_4(\Dq))^*=P_4(-\Dq)\circ q$. From $\Dq(qf)=qf+q\Dq f$, i.e.\ $\Dq\circ q=q\circ(\Dq+1)$, we get $P_4(-\Dq)\circ q=q\circ P_4(-\Dq-1)$, and
\[
P_4(-\Dq-1)=\prod_{k=1}^4\Bigl(-\Bigl(\Dq+\frac{5-k}5\Bigr)\Bigr)=\prod_{k=1}^4\Bigl(\Dq+\frac k5\Bigr)=P_4(\Dq),
\]
because there are four sign changes and $k\mapsto5-k$ permutes the factors. Hence $\PFfour^*=\PFfour$. Equivalently: expanding $\sum_k(-\Dq)^k\circ A_k$ shows that an operator of order $4$ is self-adjoint if and only if $A_3=2\Dq A_4$ and $A_1=\Dq A_2-\Dq^3A_4$; here $2\Dq A_4=-2(1-u)=A_3$ and $\Dq A_2-\Dq^3A_4=-\frac75(1-u)+(1-u)=A_1$.
\end{proof}

\begin{lem}[The concomitant]
\label{lem:gau-conc}
\begin{enumerate}
\item For every $\cP=\sum_{k\leq4}a_k\Dq^k$ with coefficients in a differential ring $\cR_0$ (for instance $\bbK$ or $\bbQ[u]$) and all $f,g$ in a differential ring extending $\cR_0$: $g\,\cP f-f\,\cP^*g=\Dq\,\Conc_\cP(f,g)$.
\item For $f,g\in\mathcal S$, $\Conc(f,g)$ lies in $\bbQ$, and $\Conc$ is a $\bbQ$-bilinear alternating form on $\mathcal S$.
\item The jet matrix is $\Jet_{ab}=\sum_{k=a+b+1}^4(-1)^{k-1-a}\binom{k-1-a}{b}\Dq^{k-1-a-b}A_k$, explicitly
\[
\Jet=\begin{pmatrix}0&\frac25(1-u)&1-u&-u\\-\frac25(1-u)&0&u&0\\-(1-u)&-u&0&0\\u&0&0&0\end{pmatrix};
\]
so $\Jet_{ab}=0$ for $a+b\geq4$, $\Jet_{3-j,j}=(-1)^ju$, and $\det\Jet=u^4$.
\item For every basis $\mathbf x$ of $\mathcal S$, $\Conc_{\mathbf x}=\Wr_{\mathbf x}^T\Jet\,\Wr_{\mathbf x}$. In the standard basis
\[
\Conc_{\mathbf f}=\begin{pmatrix}0&-\frac75&2&\ast\\\frac75&0&0&-1\\-2&0&0&0\\-\ast&1&0&0\end{pmatrix},\qquad \operatorname{Pf}(\Conc_{\mathbf f})=2,\qquad\det\Conc_{\mathbf f}=4,
\]
where $\ast=\Conc(f_0,y_3)$ is not needed. Hence $\Conc$ is non-degenerate on $\mathcal S$, $\det\Wr_{\mathbf x}\neq0$ for every basis $\mathbf x$, and $u^2\det\Wr_{\mathbf f}\in\{2,-2\}$.
\end{enumerate}
\end{lem}

\begin{proof}
(i) Fix $k$, and put $w:=a_kg$ and, for $0\leq j\leq k$, $\mathsf t_j:=(-1)^j\Dq^jw\,\Dq^{k-j}f$; these letters are local. Then $\Dq\bigl[(-1)^j\Dq^jw\,\Dq^{k-1-j}f\bigr]=\mathsf t_j-\mathsf t_{j+1}$, and the sum over $j=0,\dots,k-1$ telescopes to $\mathsf t_0-\mathsf t_k=w\,\Dq^kf-f\,(-\Dq)^kw=g\,a_k\Dq^kf-f\,(-\Dq)^k(a_kg)$. Summing over $k$ gives (i); this is an identity of differential polynomials in $f$ and $g$.

(ii) For $f,g\in\mathcal S$, (i) and \cref{lem:gau-adj} give $\Dq\Conc(f,g)=g\,\PFfour f-f\,\PFfour g=0$; as $\Conc(f,g)\in\bbK[\ell]$, \cref{lem:gau-const} gives $\Conc(f,g)\in\bbQ$. Bilinearity is clear. For alternation consider $\Conc(f,g)+\Conc(g,f)$ as a differential polynomial $Q=\sum_{a,b}q_{ab}f_ag_b$ in two indeterminates, $f_a=\Dq^af$, $g_b=\Dq^bg$, $q_{ab}\in\bbK$. By (i) and \cref{lem:gau-adj}, $\Dq Q=(g\,\PFfour f-f\,\PFfour g)+(f\,\PFfour g-g\,\PFfour f)=0$ identically. Suppose $Q\neq0$ and let $s$ be the largest value of $a+b$ with $q_{ab}\neq0$. In $\Dq Q$ the coefficient of $f_0g_{s+1}$ is $q_{0s}$, and for $c\geq1$, $c+d=s+1$, the coefficient of $f_cg_d$ is $q_{c-1,d}+q_{c,d-1}$ (with $q_{c,-1}:=0$). These vanish, so successively $q_{0,s}=q_{1,s-1}=\dots=q_{s,0}=0$, a contradiction. Hence $Q=0$ and $\Conc(f,g)=-\Conc(g,f)$.

(iii) Expanding $\Dq^j(A_kg)=\sum_b\binom jb\Dq^{j-b}A_k\,\Dq^bg$ and putting $a=k-1-j$ gives the formula; the entries follow from it with $\Dq(1-u)=1-u$, $\Dq u=-(1-u)$. For example $\Jet_{01}=-A_2+2\Dq A_3-3\Dq^2A_4=(\frac75-4+3)(1-u)$, $\Jet_{02}=A_3-3\Dq A_4=1-u$ and $\Jet_{00}=A_1-\Dq A_2+\Dq^2A_3-\Dq^3A_4=0$. Only $k=4$ contributes when $a+b=3$, giving $(-1)^{3-a}A_4=(-1)^bu$, and no $k\leq4$ contributes when $a+b\geq4$. So $\Jet$ vanishes below its antidiagonal, whose entries are $u,-u,u,-u$; the reversal of four letters is an even permutation, and $\det\Jet=u^4$.

(iv) The identity is (iii) evaluated on the basis: $(\Wr_{\mathbf x}^T\Jet\Wr_{\mathbf x})_{ij}=\sum_{a,b}\Dq^ax_i\,\Jet_{ab}\,\Dq^bx_j=\Conc(x_i,x_j)$. By (ii) every entry of $\Conc_{\mathbf f}$ is a rational number, which may be read off at $u=0$.

\emph{Entries between power series.} At $u=0$ the matrix $\Jet$ is $\frac25(E_{01}-E_{10})+E_{02}-E_{20}$, so for power-series solutions $\Conc(f,g)=\bigl[\frac25(f\Dq g-g\Dq f)+(f\Dq^2g-g\Dq^2f)\bigr](0)$. For $f=u^k+O(u^3)$ one has $(f,\Dq f,\Dq^2f)(0)=(\delta_{k0},-\delta_{k1},-\delta_{k1}+2\delta_{k2})$, i.e.\ $(1,0,0)$, $(0,-1,-1)$, $(0,0,2)$ for $f_0,f_1,f_2$. Hence $\Conc(f_0,f_1)=-\frac25-1=-\frac75$, $\Conc(f_0,f_2)=2$ and $\Conc(f_1,f_2)=0$.

\emph{Entries with $y_3$.} Write $\Dq^by_3=(\Dq^b\varpi_1)\ell+Y_b$ with $Y_b\in\bbK$. Then $\Conc(f,y_3)=\Conc(f,\varpi_1)\,\ell+\sum_{a,b}\Dq^af\,\Jet_{ab}\,Y_b$, and since $\Conc(f,y_3)$ is a constant, it equals the coefficient of $u^0$ in $\sum_{a,b}\Dq^af\,\Jet_{ab}Y_b$. With $\Dq\ell=1-u^{-1}$, $\Dq^2\ell=u^{-1}-u^{-2}$, $\Dq^3\ell=-2u^{-3}+3u^{-2}-u^{-1}$,
\begin{align*}
Y_0&=\rho_3, & Y_1&=\varpi_1\Dq\ell+\Dq\rho_3,\\
Y_2&=2\Dq\varpi_1\Dq\ell+\varpi_1\Dq^2\ell+\Dq^2\rho_3, & Y_3&=3\Dq^2\varpi_1\Dq\ell+3\Dq\varpi_1\Dq^2\ell+\varpi_1\Dq^3\ell+\Dq^3\rho_3,
\end{align*}
and $\varpi_1=u+O(u^2)$, $\Dq\varpi_1=-1+O(u)$, $\rho_3\in u^3\bbQ\fps u$ give $Y_0\in u^3\bbQ\fps u$, $Y_1=-1+O(u)$, $Y_2=u^{-1}+O(1)$ and $Y_3=u^{-2}+O(u^{-1})$. Ordering by the derivatives of $f$,
\begin{multline*}
\sum_{a,b}\Dq^af\,\Jet_{ab}Y_b=f\bigl[\tfrac25(1-u)Y_1+(1-u)Y_2-uY_3\bigr]+\Dq f\bigl[-\tfrac25(1-u)Y_0+uY_2\bigr]\\
+\Dq^2f\bigl[-(1-u)Y_0-uY_1\bigr]+\Dq^3f\,uY_0 .
\end{multline*}
The first bracket is a power series (its terms in $u^{-1}$ cancel), the second is $1+O(u)$ and the third and fourth are $O(u)$. Hence $\Conc(f,y_3)=(\Dq f)(0)$ for every power series $f$ with $f(0)=0$, so $\Conc(f_1,y_3)=-1$ and $\Conc(f_2,y_3)=0$.

For an alternating $4\times4$ matrix $(c_{ij})$, $\operatorname{Pf}=c_{01}c_{23}-c_{02}c_{13}+c_{03}c_{12}$, here $0-2\cdot(-1)+\ast\cdot0=2$, and $\det=\operatorname{Pf}^2=4$. So $\Conc$ is non-degenerate, and from $\det\Conc_{\mathbf x}=(\det\Wr_{\mathbf x})^2\det\Jet=(\det\Wr_{\mathbf x})^2u^4$ we get $\det\Wr_{\mathbf x}\neq0$ for every basis. Finally, differentiating the rows of $\Wr_{\mathbf f}$ and using $\Dq^4x=-A_4^{-1}\sum_{k\leq3}A_k\Dq^kx$ for $x\in\mathcal S$ gives Abel's identity $\Dq\det\Wr_{\mathbf f}=-(A_3/A_4)\det\Wr_{\mathbf f}=2u^{-1}(1-u)\det\Wr_{\mathbf f}$; with $\Dq(u^2)=-2u(1-u)$ this is $\Dq(u^2\det\Wr_{\mathbf f})=0$, so $u^2\det\Wr_{\mathbf f}\in\bbQ$ by \cref{lem:gau-const}, and its square is $\det\Conc_{\mathbf f}=4$.
\end{proof}

% Section 3, part 2 of 7: transversality, isotropy, the conifold expansion, special geometry, (CYC), the cyclic system.

\begin{lem}[Transversality]
\label{lem:gau-transv}
For every basis $\mathbf x$ of $\mathcal S$ and $a,b\in\{0,\dots,3\}$, in the fraction field of $\bbK[\ell]$,
\[
\sum_{i,j}(\Conc_{\mathbf x}^{-1})_{ij}\,\Dq^ax_i\,\Dq^bx_j=(\Jet^{-1})_{ab},
\]
where $(\Jet^{-1})_{ab}=0$ for $a+b\leq2$ and $(\Jet^{-1})_{03}=u^{-1}$. In particular the transversality sums of \cref{def:gau-conc} do not depend on the basis, $\Upsilon_0=\Upsilon_1=\Upsilon_2=0$ and $\Upsilon_3=u^{-1}$.
\end{lem}

\begin{proof}
$\bbK[\ell]$ is a domain and $\det\Wr_{\mathbf x}\neq0$ (\cref{lem:gau-conc}~(iv)), so $\Conc_{\mathbf x}=\Wr_{\mathbf x}^T\Jet\Wr_{\mathbf x}$ gives $\Jet^{-1}=\Wr_{\mathbf x}\Conc_{\mathbf x}^{-1}\Wr_{\mathbf x}^T$ over the fraction field; its entry $(a,b)$ is the left side. Let $J$ be the $4\times4$ exchange matrix, $J_{ac}=\delta_{a+c,3}$. The entry $(\Jet J)_{ac}=\Jet_{a,3-c}$ vanishes for $c<a$, because then $a+3-c\geq4$ (\cref{lem:gau-conc}~(iii)); so $\Jet J$ is upper triangular with diagonal entries $\Jet_{c,3-c}=(-1)^{3-c}u$, and so is its inverse $J\Jet^{-1}$, with diagonal entries $1/\Jet_{c,3-c}$. The entry $(J\Jet^{-1})_{ca}=(\Jet^{-1})_{3-c,a}$ vanishes for $a<c$, i.e.\ $(\Jet^{-1})_{ba}=0$ for $a+b<3$, and $(\Jet^{-1})_{03}=(J\Jet^{-1})_{33}=1/\Jet_{30}=u^{-1}$. The sums $\Upsilon_k$ are the entries $(0,k)$, and $\widetilde\Upsilon_k$ is built from the bivector $\sum_{i,j}(\Conc_{\mathbf x}^{-1})_{ij}x_i\otimes x_j\in\mathcal S\otimes\mathcal S$, which does not depend on $\mathbf x$: for the basis $\mathbf xM$, $M\in GL_4(\bbQ)$, the Gram matrix is $M^T\Conc_{\mathbf x}M$.
\end{proof}

\begin{lem}[Isotropy, formal Picard--Lefschetz]
\label{lem:gau-iso}
\begin{enumerate}
\item $\sh$ commutes with $\Dq$, fixes $f_0,f_1,f_2$ and maps $y_3$ to $y_3+\varpi_1$; so $\sh(\mathcal S)=\mathcal S$.
\item $\Conc(f,\varpi_1)=0$ for every $f\in\mathcal S\cap\bbK$, and $\Conc(y_3,\varpi_1)=1$.
\item On $\mathcal S\cap\bbK$ the form $\Conc$ has rank $2$ and radical $\bbQ\varpi_1$, and a plane in $\mathcal S\cap\bbK$ is isotropic if and only if it contains $\varpi_1$. In particular $\operatorname{span}(y_0,\varpi_1)$ is isotropic for every power-series solution $y_0$; and $\Conc(f_0,f_1+\gamma'f_2)=-\frac75+2\gamma'$ vanishes only for $\gamma'=\gamma=\frac7{10}$.
\end{enumerate}
\end{lem}

\begin{proof}
(i) $\sh$ is an automorphism with inverse $\ell\mapsto\ell-1$. The maps $\Dq\circ\sh$ and $\sh\circ\Dq$ are additive and both satisfy $X(fg)=X(f)\sh(g)+\sh(f)X(g)$; they agree on $\bbK$, where $\sh$ is the identity, and on $\ell$, since $\Dq(\ell+1)=\Dq\ell=\sh(\Dq\ell)$ because $\Dq\ell\in\bbK$. By the twisted Leibniz rule they agree on every monomial $a\ell^k$, hence on $\bbK[\ell]$. Clearly $\sh f_k=f_k$, and $\sh y_3=\varpi_1(\ell+1)+\rho_3=y_3+\varpi_1$.

(ii) $\Conc(f,g)$ is a polynomial in the derivatives of $f$ and $g$ with coefficients in $\bbK$, so $\sh\Conc(f,g)=\Conc(\sh f,\sh g)$ by (i). For $f\in\mathcal S\cap\bbK$ the number $\Conc(f,y_3)\in\bbQ$ (\cref{lem:gau-conc}~(ii)) is fixed by $\sh$, so $\Conc(f,y_3)=\Conc(f,y_3+\varpi_1)=\Conc(f,y_3)+\Conc(f,\varpi_1)$, i.e.\ $\Conc(f,\varpi_1)=0$. By alternation and \cref{lem:gau-conc}~(iv), $\Conc(y_3,\varpi_1)=-\Conc(f_1,y_3)-\gamma\Conc(f_2,y_3)=1$.

(iii) By \cref{lem:gau-conc}~(iv) the Gram matrix of $\Conc$ on $(f_0,f_1,f_2)$ is $\frac75(E_{10}-E_{01})+2(E_{02}-E_{20})$, of rank $2$, with kernel spanned by $(0,1,\gamma)$, i.e.\ by $\varpi_1$; this also gives the last assertion. If a plane $\Lambda\subset\mathcal S\cap\bbK$ does not contain $\varpi_1$, then $\mathcal S\cap\bbK=\Lambda\oplus\bbQ\varpi_1$ with $\varpi_1$ in the radical, so $\Conc|_\Lambda$ has rank $2$ and $\Lambda$ is not isotropic. A plane containing $\varpi_1$ is spanned by $\varpi_1$ and some $y_0$, and it is isotropic by (ii) and alternation.
\end{proof}

\begin{lem}[The conifold expansion is a ring map]
\label{lem:gau-confexp}
Let $(\alpha,\beta)\in\bbQ^2$, $\varpi_0=f_0+\alpha\varpi_1+\beta f_2$ and $\tc=\varpi_1/\varpi_0$ (\cref{def:set-frames}), and extend $\Dq$ from $\bbK$ to $\bbQ\lser s$, $s^5=u$, by $\Dq s=-(1-s^5)/(5s^4)$. The assignment of \cref{def:set-confexp} on the generators of $\Nring$ extends to a unique ring homomorphism $\Cexp_{\alpha,\beta}\colon\Nring\to\bbQ\lser s$, and $\Cexp_{\alpha,\beta}\circ\Dq=\Dq\circ\Cexp_{\alpha,\beta}$. For the elements $C_0=c_0$, $I_1=c_0\logq+i_0$, $T=I_1/C_0$ and $C_1=\mathsf W/c_0^2$ of $\Nring$ (\cref{def:set-confexp}), which evaluate to the functions of \cref{def:set-genus0},
\[
\Cexp_{\alpha,\beta}(C_0)=\varpi_0,\qquad \Cexp_{\alpha,\beta}(I_1)=\varpi_1,\qquad \Cexp_{\alpha,\beta}(\mathsf W)=\varpi_0\Dq\varpi_1-\varpi_1\Dq\varpi_0=\varpi_0^2\,\Dq\tc ,
\]
so the conifold expansion sends $C_0$, $I_1$, $T$ and $C_1$ to $\varpi_0$, $\varpi_1$, $\tc$ and $\Dq\tc$.
\end{lem}

\begin{proof}
Write $\Cexp:=\Cexp_{\alpha,\beta}$, and recall $\Xi=-\log(1-u)\in\bbQ\fps u$ (\cref{def:set-confexp}), so that $\Dq\Xi=-1$ and $\Dq^k\Xi=0$ for $k\geq2$. By \cref{def:set-confexp}, $\Cexp(L)=s^{-1}$, $\Cexp(\logq)=-\Xi$, $\Cexp(c_j)=\Dq^j\varpi_0$ and $\Cexp(i_j)=\Dq^j\varpi_1-j\Dq^{j-1}\varpi_0+\Xi\Dq^j\varpi_0=\Dq^j(\varpi_1+\Xi\varpi_0)$, the last equality by the Leibniz rule. $\Nring$ is the localisation of a polynomial ring at $L$, $c_0$ and $\mathsf W$, so the assignment extends to a unique ring homomorphism as soon as these three go to units. $s^{-1}$ and $\varpi_0$ ($\varpi_0(0)=1$) are units, and
\[
\Cexp(\mathsf W)=\varpi_0^2+\varpi_0\bigl(\Dq\varpi_1-\varpi_0+\Xi\Dq\varpi_0\bigr)-\Dq\varpi_0\bigl(\varpi_1+\Xi\varpi_0\bigr)=\varpi_0\Dq\varpi_1-\varpi_1\Dq\varpi_0=\varpi_0^2\Dq\tc,
\]
a unit because $\Dq\tc(0)=-\tc'(0)=-1$ (\cref{def:set-frames}). Also $\Cexp(I_1)=-\Xi\varpi_0+\varpi_1+\Xi\varpi_0=\varpi_1$.

$\Dq\circ\Cexp$ and $\Cexp\circ\Dq$ are additive maps $\delta$ with $\delta(xy)=\delta(x)\Cexp(y)+\Cexp(x)\delta(y)$, so it suffices to compare them on the generators of the polynomial ring: on the inverted elements $L$, $c_0$, $\mathsf W$ such a map satisfies $\delta(x^{-1})=-\Cexp(x)^{-2}\delta(x)$. For $L$: $\Dq(s^{-1})=(s^{-6}-s^{-1})/5=\Cexp((L^6-L)/5)$; for $\logq$: $\Dq(-\Xi)=1$; for $c_0,c_1,c_2,i_0,i_1$ the images $\Dq^j\varpi_0$ and $\Dq^j(\varpi_1+\Xi\varpi_0)$ are shifted in $j$ by $\Dq$. In $\Nring$, $\Dq c_3$ is fixed by $\PFfour c_0=0$ (coefficients written through $u=L^{-5}$), and $\PFfour\varpi_0=0$ gives $\Cexp(\Dq c_3)=\Dq^4\varpi_0$. Finally $\Dq i_2$ is fixed by $\Conc(c_0,I_1)=0$, where $\Dq^3I_1=c_3\logq+3c_2+\Dq i_2$ and $\Conc$ has its coefficients in $\bbQ[u]$, $u=L^{-5}$ (\cref{def:gau-conc}); this relation is affine in $\Dq i_2$ with the unit coefficient $\Jet_{03}c_0=-L^{-5}c_0$, and by the cases treated $\Cexp(\Dq^bI_1)=\Dq^b\varpi_1$ for $b\leq2$. Its image under $\Cexp$ differs from $\Conc(\varpi_0,\varpi_1)=0$ (\cref{lem:gau-iso}~(ii)) only in the last slot, so $-u\varpi_0\bigl(\Cexp(\Dq^3I_1)-\Dq^3\varpi_1\bigr)=0$, and $\Cexp(\Dq i_2)=\Dq^3\varpi_1+\Xi\Dq^3\varpi_0-3\Dq^2\varpi_0=\Dq^3(\varpi_1+\Xi\varpi_0)=\Dq\Cexp(i_2)$. The last assertion follows from $\Cexp(I_1/c_0)=\tc$ and $\Cexp(\mathsf W/c_0^2)=\Dq\tc$.
\end{proof}

\begin{lem}[Special geometry]
\label{lem:gau-sg}
Let $y_0\in\bbK^\times$ and $t\in\bbK$ with $\Dq t\neq0$.
\begin{enumerate}
\item[(a)] For every basis $\mathbf x$ of $\mathcal S$: $\widetilde\Upsilon_1=\widetilde\Upsilon_2=0$ and $\widetilde\Upsilon_3=1/(u\,y_0^2(\Dq t)^3)$.
\item[(b)] Suppose moreover that $y_0$ and $y_1:=ty_0$ lie in $\mathcal S\cap\bbK$ and span an isotropic plane; for instance $y_1=\varpi_1$ and $y_0$ any power-series solution not proportional to $\varpi_1$ (\cref{lem:gau-iso}). Then $\mathcal S$ has a basis $(x_0,\dots,x_3)$ with $x_0=y_0$, $x_1=y_1$, $\Conc(x_0,x_3)=\Conc(x_1,x_2)=1$ and $\Conc(x_0,x_2)=\Conc(x_1,x_3)=\Conc(x_2,x_3)=0$. With $\phi_i:=x_i/y_0$ (so $\phi_0=1$, $\phi_1=t$) and $\mathrm{Yuk}:=\de_t^2\phi_2$:
\begin{enumerate}
\item $\de_t^2\phi_3=-t\,\mathrm{Yuk}$; hence $\de_t^2(y/y_0)\in\mathrm{Yuk}\cdot\operatorname{span}_\bbQ(1,t)$ for every $y\in\mathcal S$;
\item $\mathrm{Yuk}=1/(u\,y_0^2(\Dq t)^3)$.
\end{enumerate}
\end{enumerate}
\end{lem}

\begin{proof}
(a) Put $\dot t:=\Dq t$, a unit of $\bbK$, and (local notation) $\Upsilon^{\Dq}_m:=\sum_{i,j}(\Conc_{\mathbf x}^{-1})_{ij}(x_i/y_0)\Dq^m(x_j/y_0)$; $\Upsilon^{\Dq}_0=0$ because $\Conc_{\mathbf x}^{-1}$ is antisymmetric. The Leibniz rule for $x_j=y_0\cdot(x_j/y_0)$ gives $\Upsilon_k=\sum_m\binom km\,y_0\Dq^{k-m}y_0\,\Upsilon^{\Dq}_m$, so
\[
\Upsilon_1=y_0^2\Upsilon^{\Dq}_1,\quad \Upsilon_2=2y_0\Dq y_0\,\Upsilon^{\Dq}_1+y_0^2\Upsilon^{\Dq}_2,\quad \Upsilon_3=3y_0\Dq^2y_0\,\Upsilon^{\Dq}_1+3y_0\Dq y_0\,\Upsilon^{\Dq}_2+y_0^2\Upsilon^{\Dq}_3 .
\]
As operators $\Dq=\dot t\de_t$, $\Dq^2=(\Dq\dot t)\de_t+\dot t^2\de_t^2$ and $\Dq^3=(\Dq^2\dot t)\de_t+3\dot t(\Dq\dot t)\de_t^2+\dot t^3\de_t^3$, hence $\Upsilon^{\Dq}_1=\dot t\,\widetilde\Upsilon_1$, $\Upsilon^{\Dq}_2=(\Dq\dot t)\widetilde\Upsilon_1+\dot t^2\widetilde\Upsilon_2$ and $\Upsilon^{\Dq}_3=(\Dq^2\dot t)\widetilde\Upsilon_1+3\dot t(\Dq\dot t)\widetilde\Upsilon_2+\dot t^3\widetilde\Upsilon_3$. By \cref{lem:gau-transv}, $\Upsilon_1=\Upsilon_2=0$ and $\Upsilon_3=u^{-1}$; since $y_0$ and $\dot t$ are units, solving successively gives $\Upsilon^{\Dq}_1=\Upsilon^{\Dq}_2=0$, $\Upsilon^{\Dq}_3=1/(uy_0^2)$, and then (a). Isotropy is not used in (a).

(b) $y_0$ and $y_1$ are linearly independent because $t$ is not constant, so $\Lambda:=\operatorname{span}(y_0,y_1)$ is an isotropic plane. As $\Conc$ is non-degenerate on the four-dimensional space $\mathcal S$ (\cref{lem:gau-conc}~(iv)), $\Lambda^\perp$ is two-dimensional and contains $\Lambda$, so $\Lambda^\perp=\Lambda$, and the map $\mathcal S\to\Lambda^*$, $x\mapsto\Conc(\cdot,x)|_\Lambda$, is onto. Choose $x_3$ with $\Conc(y_0,x_3)=1$, $\Conc(y_1,x_3)=0$, and $x_2'$ with $\Conc(y_0,x_2')=0$, $\Conc(y_1,x_2')=1$, and put $x_2:=x_2'-\Conc(x_2',x_3)\,y_0$; then $\Conc(y_0,x_2)=0$, $\Conc(y_1,x_2)=1$ and $\Conc(x_2,x_3)=0$. The Gram matrix of $(x_0,\dots,x_3)$ is $E_{03}+E_{12}-E_{21}-E_{30}$, invertible, so this is a basis, and $\Conc_{\mathbf x}^{-1}=-\Conc_{\mathbf x}$. In this basis the definition of $\widetilde\Upsilon_k$ reads
\[
\widetilde\Upsilon_k=-\phi_0\,\de_t^k\phi_3+\phi_3\,\de_t^k\phi_0-\phi_1\,\de_t^k\phi_2+\phi_2\,\de_t^k\phi_1 .
\]
For $k\geq2$, $\de_t^k1=\de_t^kt=0$, so $\widetilde\Upsilon_k=-\de_t^k\phi_3-t\,\de_t^k\phi_2$. For $k=2$ this gives $0=-\de_t^2\phi_3-t\,\mathrm{Yuk}$, which is (i); for $k=3$ it gives $\widetilde\Upsilon_3=\de_t(t\,\mathrm{Yuk})-t\,\de_t\mathrm{Yuk}=\mathrm{Yuk}$, which is (ii) by (a). Finally, for $y=\sum_ic_ix_i$ with $c_i\in\bbQ$, $\de_t^2(y/y_0)=c_2\mathrm{Yuk}-c_3t\,\mathrm{Yuk}$.
\end{proof}

In a conifold frame, $\mathrm{Yuk}=1/(u\varpi_0^2(\Dq\tc)^3)$ (with $u=A_4$) is the conifold expansion of $\kappa^{\mathrm{LR}}_1$ of \cref{prop:set-cyclicLR}; it is normalised by the Darboux form of (b), the three-point function $\eta(H*H,H)=5\kappa^{\mathrm{LR}}_1$ carries the factor $5$ of $\eta$, and the coefficient $\kappa_1=\mathrm{Yuk}\cdot\Dq\tc$ of \cref{prop:gau-quintic} carries none.

\begin{thm}[The cyclic form of $\PFfour$, (CYC)]
\label{thm:gau-cyc}
Let $y_0\in\mathcal S\cap\bbK$ be non-zero and not proportional to $\varpi_1$, and let $\varpi_1$ be taken at any non-zero scale. Then, as operators in $\bbK[\Dq]$,
\begin{equation}
\label{eq:gau-cyc}
\PFfour=\Mcyc(y_0,\varpi_1)=\frac1{y_0}\circ\Dq\circ\frac1{\Dq t}\circ\Dq\circ\bigl(u\,y_0^2(\Dq t)^2\bigr)\circ\Dq\circ\frac1{\Dq t}\circ\Dq\circ\frac1{y_0},\qquad t=\frac{\varpi_1}{y_0}.
\end{equation}
More generally $\PFfour=\Mcyc(y_0,y_1)$ for every basis $(y_0,y_1)$ of an isotropic plane in $\mathcal S\cap\bbK$; these are the planes that contain $\varpi_1$ (\cref{lem:gau-iso}~(iii)). $\Mcyc(y_0,y_1)$ is unchanged under $y_0\mapsto\mu y_0$, $y_1\mapsto\mu'y_1$ ($\mu,\mu'\in\bbQ^\times$) and $y_1\mapsto y_1+ay_0$ ($a\in\bbQ$). In particular $\PFfour=\Mcyc(\varpi_0,\varpi_1)$ in every conifold frame of \cref{def:set-frames}.
\end{thm}

\begin{proof}
Put $t:=y_1/y_0$ and $\dot t:=\Dq t$. As operators $\Dq\circ\dot t^{-1}\circ\Dq=\dot t\,\de_t^2$, so
\[
\Mcyc(y_0,y_1)=\frac{\dot t}{y_0}\circ\de_t^2\circ\bigl(u\,y_0^2\dot t^3\bigr)\circ\de_t^2\circ\frac1{y_0}.
\]
For $y\in\mathcal S$, \cref{lem:gau-sg}~(b) gives $u\,y_0^2\dot t^3\de_t^2(y/y_0)=c_2-c_3t$ with constants $c_2,c_3$, which $\de_t^2$ kills; so $\Mcyc(y_0,y_1)$ annihilates $\mathcal S$. The coefficient of $\Dq^4$ in \eqref{eq:gau-Mcyc} is the product of the five multiplication factors,
\[
y_0^{-1}\cdot\dot t^{-1}\cdot uy_0^2\dot t^2\cdot\dot t^{-1}\cdot y_0^{-1}=u=A_4.
\] Hence $E:=\Mcyc(y_0,y_1)-\PFfour=\sum_{k\leq3}e_k\Dq^k\in\bbK[\Dq]$ annihilates $\mathcal S$, i.e.\ the row vector $(e_0,\dots,e_3)$ satisfies $(e_0,\dots,e_3)\Wr_{\mathbf f}=0$; as $\det\Wr_{\mathbf f}\neq0$ (\cref{lem:gau-conc}~(iv)), $E=0$. Every step takes place in $\bbK$ or $\bbK[\ell]$, where $y_0$ and $\dot t$ are units; the condition $y_0(0)\neq0$ is not used. Under $y_0\mapsto\mu y_0$, $y_1\mapsto\mu'y_1$ the multiplication factors of \eqref{eq:gau-Mcyc} are multiplied by $\mu^{-1},\mu\mu'^{-1},\mu'^2,\mu\mu'^{-1},\mu^{-1}$, with product $1$, and constants commute with $\Dq$; under $y_1\mapsto y_1+ay_0$, $t\mapsto t+a$ and $\dot t$ is unchanged. The plane of the first statement contains $\varpi_1$ and is isotropic by \cref{lem:gau-iso}; in a conifold frame $\varpi_0$ is a power-series solution with $\varpi_0(0)=1$, not proportional to $\varpi_1$.
\end{proof}

\begin{cor}[The cyclic system is $\PFfive$]
\label{cor:gau-cycsys}
Let $(y_0,y_1)$ be as in \cref{thm:gau-cyc}, $t=y_1/y_0$, and
\[
\kappa=(\kappa_0,\dots,\kappa_4):=\Bigl(\Dq t,\ \frac{L^5}{y_0^2(\Dq t)^2},\ \Dq t,\ y_0,\ y_0\Bigr),\qquad L^5:=u^{-1},
\]
units of $\bbK$. Let $(\cR,\Dq)$ be a differential ring extending $(\bbK,\Dq)$ and $z\in\cR$ with $\Dq z=0$.
\begin{enumerate}
\item For $Y_0,\dots,Y_4\in\cR$ the cyclic system
\begin{equation}
\label{eq:gau-system}
z\Dq Y_a=\kappa_aY_{a+1}\qquad(a\in\bbZ/5)
\end{equation}
holds if and only if $y_0Y_0$ solves $\PFfive$ and $Y_{a+1}=z\Dq Y_a/\kappa_a$ for $a=0,\dots,3$.
\item If $z$ is invertible in $\cR$, then for every $w\in\cR$ the same holds with $z\Dq$ replaced by $z\Dq+w$ in \eqref{eq:gau-system} and in $\PFfive$.
\end{enumerate}
\end{cor}

\begin{proof}
Inserting the $\kappa_a$ and comparing with \eqref{eq:gau-Mcyc} gives the identity of operators
\[
\kappa_4^{-1}\Dq\,\kappa_3^{-1}\Dq\,\kappa_2^{-1}\Dq\,\kappa_1^{-1}\Dq\,\kappa_0^{-1}\Dq\circ y_0^{-1}=y_0^{-1}\Dq\,y_0^{-1}\Dq\,\dot t^{-1}\Dq\,(uy_0^2\dot t^2)\Dq\,\dot t^{-1}\Dq\,y_0^{-1}=y_0^{-1}\Dq\circ\Mcyc(y_0,y_1),
\]
with $\dot t=\Dq t$ and $\kappa_1^{-1}=uy_0^2\dot t^2$. (i) Given the four relations $Y_{a+1}=z\Dq Y_a/\kappa_a$, $a=0,\dots,3$, the closing relation $z\Dq Y_4=\kappa_4Y_0$ reads $Y_0=z^5\kappa_4^{-1}\Dq\kappa_3^{-1}\Dq\cdots\kappa_0^{-1}\Dq\,Y_0$, because $z$ commutes with $\Dq$; by the identity and \cref{thm:gau-cyc} this is $y_0Y_0=z^5\Dq\,\PFfour(y_0Y_0)$, which is $\PFfive(y_0Y_0)=0$ by \cref{lem:set-pf5}. (That lemma is an identity of operators with coefficients in $\bbQ[q][z]$: its proof uses only $\Dq\circ q=q\circ(\Dq+1)$, which holds in $\cR$ for $q=(1-u)/5^5\in\bbK$.) The same chain of equivalences read in the other direction gives the converse. (ii) If $z$ is invertible, the map of the ring of differential operators $\cR[\Dq]$ that fixes $\cR$ and sends $\Dq$ to $\Dq+w/z$ is a ring automorphism, since $(\Dq+w/z)\circ f-f\circ(\Dq+w/z)=\Dq(f)$ for $f\in\cR$; it maps $z\Dq$ to $z\Dq+w$, and applying it to the operator identities used in (i) gives (ii).
\end{proof}

The corollary is applied in $\bbQ(\zeta)\lser s\lser z$ with $w=\zeta^iL=\zeta^is^{-1}$, to amplitudes of formal WKB solutions with exponential $\re^{u_i/z}$, $\Dq u_i=\zeta^iL$: the expanded rows (\cref{prop:gau-quintic}) and the canonical gauge applied to the model columns (\cref{thm:gau-R0}), and, through $z=s^4\zh$, to the hatted series of \cref{lem:gau-HPF}; with $w=0$ it applies in $\bbQ\lser z\lser u$ to solutions without exponential, such as the vanishing solutions of \cref{def:gau-van}. In a conifold frame $\kappa$ is the vector of \cref{prop:gau-quintic}. It does not assert that WKB solutions with rows free of logarithms exist; for the rows of \cref{def:set-rows} this is part of \cref{thm:gau-R0}.

% Section 3, part 3 of 7: the quintic system and the pulled-back model; Theorem U (statement); the commutant.

%-------------------------------------------------------------------------------
\subsection{The quintic system and the pulled-back model}
\label{sec:gau-systems}

\begin{defn}[Cyclic systems with one simple pole; regular gauges]
\label{def:gau-onepole}
Let $\Bbbk$ be a field of characteristic $0$. A \emph{cyclic system with one simple pole} over $\Bbbk\fps u$ is a system $z\de_uY=AY$ with $A=\sum_{a\in\bbZ/5}b_aE_{a,a+1}$ whose coefficients have the \emph{one-pole shape}: $b_a\in\Bbbk\fps u^\times$ for $a\neq1$ and $ub_1\in\Bbbk\fps u^\times$. Its \emph{residue} is $r:=(ub_1)(0)\in\Bbbk^\times$ and its \emph{product} is $\Bprod:=\prod_ab_a\in\Bbbk\lser u$. It is \emph{self-dual} if $b_0=b_2$ and $b_3=b_4$; for $\eta=(5\,\delta_{a+b\equiv3})_{a,b}$ the entry $(c,d)$ of $\eta A$ is $5b_{d-1}\delta_{c+d\equiv4}$ and that of $A^T\eta$ is $5b_{c-1}\delta_{c+d\equiv4}$, so self-duality is $\eta A=A^T\eta$. For a second such system $\tilde A=\sum_a\tilde b_aE_{a,a+1}$, with residue $\tilde r$, the condition of \emph{equal products} is
\[
\text{(EP)}\qquad \prod_ab_a=\prod_a\tilde b_a .
\]
The \emph{gauge equation} from $\tilde A$ to $A$ is $zP'=AP-P\tilde A$ for $P=\sum_{k\geq0}z^kP_k$, $P_k\in\mathfrak{gl}_5(\Bbbk\fps u)$; then $z\de_u(P\tilde Y)=(AP-P\tilde A)\tilde Y+P\tilde A\tilde Y=AP\tilde Y$ whenever $z\de_u\tilde Y=\tilde A\tilde Y$, so $Y=P\tilde Y$ maps solutions of the second system to solutions of the first. $\Gau(A,\tilde A)$ denotes the set of all solutions $P\in\mathfrak{gl}_5(\Bbbk\fps u)\fps z$, with no condition on $P_0$; it is a $\Bbbk\fps z$-module, since $z(P\cn)'=(AP-P\tilde A)\cn$ for $\cn\in\Bbbk\fps z$. A \emph{regular gauge} is an element of $\Gau(A,\tilde A)$ with $P_0$ invertible. $P$ is \emph{graded} if $P_k=\sum_bp_{k,b}E_{b+k,b}$ for every $k$.
\end{defn}

\begin{prop}[The quintic cyclic system]
\label{prop:gau-quintic}
Let $(\alpha,\beta)\in\bbQ^2$ give a conifold frame (\cref{def:set-frames}), $\tc=\varpi_1/\varpi_0$, and
\[
\kappa:=\Bigl(\Dq\tc,\ \frac{L^5}{\varpi_0^2(\Dq\tc)^2},\ \Dq\tc,\ \varpi_0,\ \varpi_0\Bigr),\qquad L^5=u^{-1},
\]
the vector of \cref{cor:gau-cycsys} for $(y_0,y_1)=(\varpi_0,\varpi_1)$.
\begin{enumerate}
\item $\kappa_a=\Cexp_{\alpha,\beta}(C_1\kappa^{\mathrm{LR}}_a)$ for $a\in\bbZ/5$, where $C_1\kappa^{\mathrm{LR}}=(C_1,L^5/(C_0^2C_1^2),C_1,C_0,C_0)$ (\cref{prop:set-cyclicLR}) is represented in $\Nring$ by $(\mathsf W/c_0^2,\ L^5c_0^2/\mathsf W^2,\ \mathsf W/c_0^2,\ c_0,\ c_0)$. Hence the conifold expansion carries the forward relations $(z\Dq+\zeta^iL)X_a=C_1\kappa^{\mathrm{LR}}_aX_{a+1}$ ($a=0,\dots,3$) of elements of $\Nring\lser z\otimes\bbQ(\zeta)$ to the relations $(z\Dq+\zeta^iL)\Cexp(X_a)=\kappa_a\Cexp(X_{a+1})$, $L=s^{-1}$; the closing relation is then given by \cref{cor:gau-cycsys}.
\item The cyclic system \eqref{eq:gau-system} is $z\de_uY=AY$ with $A=\sum_ab_aE_{a,a+1}$, $b_a:=-\kappa_a/(1-u)$,
\[
b=\Bigl(\tc',\ -\frac1{u\varpi_0^2(1-u)^3\tc'^2},\ \tc',\ -\frac{\varpi_0}{1-u},\ -\frac{\varpi_0}{1-u}\Bigr),
\]
a self-dual cyclic system with one simple pole over $\bbQ\fps u$, with residue $r=-1$, product $\prod_ab_a=-1/(u(1-u)^5)$, and $b_a(0)=1,1,-1,-1$ for $a=0,2,3,4$.
\item More generally, for $y_0\in\mathcal S\cap\bbK$ not proportional to $\varpi_1$ and $t=\varpi_1/y_0$, the coefficients $-\kappa_a/(1-u)$ of \cref{cor:gau-cycsys} have the one-pole shape if and only if $y_0(0)\neq0$.
\end{enumerate}
\end{prop}

\begin{proof}
(i) By \cref{lem:gau-confexp}, $\Cexp(c_0)=\varpi_0$, $\Cexp(\mathsf W/c_0^2)=\Dq\tc$ and $\Cexp(L^5)=s^{-5}=u^{-1}$; hence $\Cexp(L^5c_0^2/\mathsf W^2)=u^{-1}\varpi_0^2/(\varpi_0^4(\Dq\tc)^2)=\kappa_1$, and the other four entries are immediate. Since $\Cexp$ is a ring homomorphism commuting with $\Dq$, with $\Cexp(L)=s^{-1}$, it carries each forward relation to the corresponding relation with $\kappa_a$. By \cref{cor:gau-cycsys}~(ii), with $w=\zeta^is^{-1}$, the closing relation holds if and only if $\varpi_0\Cexp(X_0)$ solves $\PFfive$ with $z\Dq$ replaced by $z\Dq+\zeta^iL$.

(ii) As $\Dq=-(1-u)\,\dd/\dd u$, the relation $z\Dq Y_a=\kappa_aY_{a+1}$ is $z\de_uY_a=b_aY_{a+1}$. With $\Dq\tc=-(1-u)\tc'$ and $L^5=u^{-1}$ the coefficients are those displayed. In a conifold frame $\varpi_0(0)=1$ and $\tc'(0)=1$ (\cref{def:set-frames}), so $b_0,b_2,b_3,b_4$ are units with the stated values and $ub_1$ is a unit with $(ub_1)(0)=-1$. Self-duality is $b_0=b_2$, $b_3=b_4$, and $\prod_ab_a=\tc'^2\cdot\bigl(-u^{-1}\varpi_0^{-2}(1-u)^{-3}\tc'^{-2}\bigr)\cdot\varpi_0^2(1-u)^{-2}=-u^{-1}(1-u)^{-5}$.

(iii) If $y_0(0)\neq0$, then $t=u/y_0(0)+O(u^2)$ since $\varpi_1=u+O(u^2)$, so $t'(0)\neq0$, and the computation of (ii) with $(\varpi_0,\tc)$ replaced by $(y_0,t)$ shows the one-pole shape. If $y_0(0)=0$, then $-\kappa_3/(1-u)=-y_0/(1-u)$ is not a unit.
\end{proof}

\begin{prop}[Equal products and $v(u)$]
\label{prop:gau-equalprod}
Let $v\in u\bbQ\fps u$ with $v'(0)\neq0$. The model system $z\,\dd Y_b/\dd v=b^{\mathrm m}_bY_{b+1}$, $b^{\mathrm m}=(1,-v^{-1},1,-1,-1)$ (\cref{prop:set-modrows}), pulled back along $v=v(u)$ is $z\de_uY=\tilde AY$ with $\tilde A=v'(u)\sum_bb^{\mathrm m}_b(v(u))E_{b,b+1}$, i.e.
\[
\tilde b=\bigl(v',\ -v'/v,\ v',\ -v',\ -v'\bigr),\qquad \prod_a\tilde b_a=-v'^5/v .
\]
\begin{enumerate}
\item If $v=u+O(u^2)$, $\tilde A$ is a self-dual cyclic system with one simple pole over $\bbQ\fps u$, with residue $\tilde r=-1$ and $\tilde b_a(0)=1,1,-1,-1$ for $a=0,2,3,4$.
\item For $v=u+O(u^2)$, the products of $\tilde A$ and of the quintic system of \cref{prop:gau-quintic} are equal if and only if $v'^5/v=1/(u(1-u)^5)$, if and only if $v=u\,\tau(u)^{-5/4}$ (\cref{def:set-tau}).
\item (Necessity of (EP).) If $A,\tilde A$ are cyclic systems with one simple pole over $\Bbbk\fps u$ and $\Gau(A,\tilde A)\neq0$, then $\prod_ab_a=\prod_a\tilde b_a$.
\end{enumerate}
In particular, for $v=v(u)$ and every conifold frame, the quintic system and the pulled-back model are two self-dual cyclic systems with one simple pole satisfying (EP), and $b_a(0)=\tilde b_a(0)$ for $a\neq1$, $r=\tilde r=-1$.
\end{prop}

\begin{proof}
The chain rule gives $z\de_u[Y(v(u))]=v'(u)\,(z\,\dd Y/\dd v)(v(u))$, hence $\tilde A$ and $\tilde b$; and $\prod_a\tilde b_a=v'^5\cdot(-v^{-1})\cdot(-1)^2=-v'^5/v$. (i) With $v=u+O(u^2)$, $v'$ is a unit with $v'(0)=1$ and $-uv'/v$ is a unit with value $-1$ at $u=0$.

(ii) The first equivalence is $\prod_ab_a=-1/(u(1-u)^5)$ (\cref{prop:gau-quintic}~(ii)). Work in $\bbQ\lser s$, $s^5=u$, where $\dd/\dd u=(5s^4)^{-1}\dd/\dd s$ has kernel $\bbQ$. Write $v=uw$ with $w\in1+u\bbQ\fps u$; an element of $1+u\bbQ\fps u$ has a unique $n$-th root in $1+u\bbQ\fps u$ for $n\geq1$ (binomial series; and $c^n=1$ for $c=1+c_ku^k+\dots$, $c_k\neq0$, is impossible, as $c^n=1+nc_ku^k+\dots$). Since $v'(0)=1$, the equation $v'^5=v\,u^{-1}(1-u)^{-5}=w(1-u)^{-5}$ is equivalent to $v'=w^{1/5}(1-u)^{-1}$. Put $v^{1/5}:=s\,w^{1/5}$ and $v^{4/5}:=s^4w^{4/5}$; then $\frac{\dd}{\dd u}v^{4/5}=\frac45v^{-1/5}v'$, and the equation is equivalent to $\frac{\dd}{\dd u}v^{4/5}=\frac45s^{-1}(1-u)^{-1}$. On the other hand $s^4/\tau(u)=\sum_{n\geq0}\frac{4}{5n+4}u^{n+4/5}$ (\cref{def:set-tau}), whose derivative is $\frac45\sum_nu^{n-1/5}=\frac45s^{-1}(1-u)^{-1}$. So the equation is equivalent to $\frac{\dd}{\dd u}\bigl(v^{4/5}-s^4/\tau\bigr)=0$, i.e.\ to $v^{4/5}-s^4/\tau\in\bbQ$; as this difference lies in $s^4\bbQ\fps u$, it is equivalent to $w^{4/5}=\tau^{-1}$, i.e.\ $w=\tau^{-5/4}$, by the uniqueness of roots. All steps are equivalences.

(iii) Let $P\in\Gau(A,\tilde A)$ be non-zero and $n$ minimal with $P_n\neq0$. The coefficient of $z^n$ in $zP'=AP-P\tilde A$ is $0=AP_n-P_n\tilde A$, hence $A^5P_n=P_n\tilde A^5$. Now $A^5=\sum_ab_ab_{a+1}\cdots b_{a+4}E_{a,a+5}=(\prod_ab_a)\,1$, and likewise for $\tilde A$, so $(\prod_ab_a-\prod_a\tilde b_a)P_n=0$ in $\mathfrak{gl}_5(\Bbbk\lser u)$, and $P_n\neq0$ gives the claim.

The last sentence combines (i), (ii) and \cref{prop:gau-quintic}~(ii).
\end{proof}

%-------------------------------------------------------------------------------
\subsection{Theorem U}
\label{sec:gau-U}

Throughout this subsection $\Bbbk$ is a field of characteristic $0$, $A=\sum_ab_aE_{a,a+1}$ and $\tilde A=\sum_a\tilde b_aE_{a,a+1}$ are cyclic systems with one simple pole over $\Bbbk\fps u$ satisfying (EP), and a prime is $\dd/\dd u$. The following notation is local to this subsection.
\begin{itemize}
\item $\mathfrak{gl}^{(m)}:=\bigoplus_a\Bbbk\lser uE_{a,a+m}$ for $m\in\bbZ/5$; an element of $\mathfrak{gl}^{(m)}$ has \emph{shift} $m$, and $\mathfrak{gl}_5(\Bbbk\lser u)=\bigoplus_m\mathfrak{gl}^{(m)}$.
\item $\Bprod^{(m)}_a:=\tilde b_a\tilde b_{a+1}\cdots\tilde b_{a+m-1}$ for $0\leq m\leq5$ (so $\Bprod^{(0)}_a=1$ and $\Bprod^{(5)}_a=\Bprod:=\prod_a\tilde b_a$), and $\mathcal Y:=\tilde A^{-1}$.
\item $\Theta_U:=\diag(d^U_0,\dots,d^U_4)$ with $d^U_0:=1$ and $d^U_{a+1}:=d^U_a\tilde b_a/b_a$; $\operatorname{dlog}\Theta_U:=\Theta_U^{-1}\Theta_U'$; $\mathfrak d:=(\det\Theta_U/\det\Theta_U(0))^{-1/5}$.
\item $\operatorname{pr}$, $\chi^U_j$, $P^\natural$ and the degree parts $P^{[\delta]}$ are defined in \cref{lem:gau-comm,lem:gau-scal}, \cref{thm:gau-U} and \cref{lem:gau-degree}. If $\Bbbk$ contains $\zeta$, put $\mathrm{gr}:=\diag(0,1,2,3,4)$, $\zeta^{i(\mathrm{gr}+1)}:=\diag(\zeta^{i(a+1)})_a$, and define the \emph{colour rotation} $\operatorname{Rot}_iP(u,z):=\zeta^{i(\mathrm{gr}+1)}P(u,\zeta^{-i}z)\zeta^{-i(\mathrm{gr}+1)}$ on matrices and $\operatorname{Rot}_iY(u,z):=\zeta^{i(\mathrm{gr}+1)}Y(u,\zeta^{-i}z)$ on vectors.
\end{itemize}
By (EP), $d^U_4\tilde b_4/b_4=\prod_a\tilde b_a/b_a=1=d^U_0$, so $d^U_{a+1}=d^U_a\tilde b_a/b_a$ holds for every $a\in\bbZ/5$. The $d^U_a$ are units of $\Bbbk\fps u$, since $\tilde b_a/b_a$ is a unit for $a\neq1$ and $\tilde b_1/b_1=(u\tilde b_1)/(ub_1)$; hence $\det\Theta_U/\det\Theta_U(0)\in1+u\Bbbk\fps u$, $\mathfrak d\in1+u\Bbbk\fps u$ is its binomial series, and $\mathfrak d'=-\frac15\tr(\operatorname{dlog}\Theta_U)\,\mathfrak d$. Entrywise, $(A\Theta_U)_{a,a+1}=b_ad^U_{a+1}=d^U_a\tilde b_a=(\Theta_U\tilde A)_{a,a+1}$, so $A\Theta_U=\Theta_U\tilde A$.

\begin{thm}[Theorem U]
\label{thm:gau-U}
Let $A$, $\tilde A$ be cyclic systems with one simple pole over $\Bbbk\fps u$ with equal products.
\begin{enumerate}
\item[(a)] (Existence.) There is $P^\natural\in\Gau(A,\tilde A)$ with $P^\natural_0=\mathfrak d\,\Theta_U$, a diagonal matrix with unit entries.
\item[(b)] (Uniqueness.) For every $P^\natural\in\Gau(A,\tilde A)$ with $P^\natural_0=\mathfrak d\,\Theta_U$, the map $\Bbbk\fps z\to\Gau(A,\tilde A)$, $\cn\mapsto P^\natural\cn(z)$, is bijective. Consequently every solution $P$ of the gauge equation with coefficients in $\mathfrak{gl}_5(\Bbbk\fps u)$ is $P^\natural\cn(z)$ for a unique $\cn\in\Bbbk\fps z$; its $z^0$-term is $P_0=\cn(0)\,\mathfrak d\,\Theta_U$, a diagonal matrix, and $P$ is a regular gauge if and only if $\cn(0)\neq0$. If $P$ is one regular gauge, the regular gauges are exactly the $P\cn(z)$ with $\cn\in\Bbbk\fps z^\times$. For the pair $(\tilde A,\tilde A)$: $\Gau(\tilde A,\tilde A)=\Bbbk\fps z\cdot1$.
\item[(c)] (Graded gauges.) Graded regular gauges exist; if $P$ is one, the graded regular gauges are exactly the $P\cn(z)$ with $\cn\in\Bbbk\fps{z^5}^\times$.
\item[(d)] (The field.) If all $b_a$, $\tilde b_a$ lie in $\bbQ\lser u$, then $P^\natural$ can be chosen in $\mathfrak{gl}_5(\bbQ\fps u)\fps z$; for every field $\Bbbk\supseteq\bbQ$ the solutions with coefficients in $\mathfrak{gl}_5(\Bbbk\fps u)$ are then $P^\natural\Bbbk\fps z$, and $P^\natural\cn$ has rational coefficients if and only if $\cn\in\bbQ\fps z$.
\end{enumerate}
\end{thm}

Part (a) is proved after \cref{lem:gau-scal}, parts (b) and (d) after \cref{lem:gau-lead}, and part (c) is \cref{prop:gau-graded}. The proof uses only the one-pole shape, (EP) and the characteristic $0$; convergence is never used.

\begin{lem}[The commutant]
\label{lem:gau-comm}
Let $F$ be a field of characteristic $0$, $\tilde b_0,\dots,\tilde b_4\in F^\times$, $\tilde A=\sum_a\tilde b_aE_{a,a+1}$, and define $\mathfrak{gl}^{(m)}$, $\Bprod^{(m)}_a$, $\Bprod$ and $\mathcal Y$ over $F$ as above.
\begin{enumerate}
\item[(a)] $\tilde A^m=\sum_a\Bprod^{(m)}_aE_{a,a+m}$ for $0\leq m\leq5$, so $\tilde A^5=\Bprod\,1$; $\mathcal Y=\sum_a\tilde b_a^{-1}E_{a+1,a}$, $\mathcal Y^j=\sum_a(\Bprod^{(j)}_a)^{-1}E_{a+j,a}$ for $0\leq j\leq5$, and $\tilde A^m=\Bprod\,\mathcal Y^{5-m}$ for $1\leq m\leq5$; $\det(x1-\tilde A)=x^5-\Bprod$.
\item[(b)] For $X=\sum_a\xi_aE_{a,a+m}\in\mathfrak{gl}^{(m)}$, $[\tilde A,X]=\sum_a(\tilde b_a\xi_{a+1}-\xi_a\tilde b_{a+m})E_{a,a+m+1}\in\mathfrak{gl}^{(m+1)}$.
\item[(c)] For every $m$, $\ker(\operatorname{ad}\tilde A)\cap\mathfrak{gl}^{(m)}=F\tilde A^m$; hence $\ker(\operatorname{ad}\tilde A)=F[\tilde A]=F[\mathcal Y]=\bigoplus_{j=0}^4F\mathcal Y^j$, of dimension $5$.
\item[(d)] For $m\in\{1,\dots,5\}$ and $X=\sum_a\xi_aE_{a,a+m}\in\mathfrak{gl}^{(m)}$ put $\lambda_m(X):=\sum_a\xi_a/\Bprod^{(m)}_a$ (a local letter). Then $\lambda_m(\tilde A^m)=5$, $\operatorname{ad}\tilde A(\mathfrak{gl}^{(m-1)})=\ker\lambda_m$, $\mathfrak{gl}^{(m)}=F\tilde A^m\oplus\ker\lambda_m$, and $\mathfrak{gl}_5(F)=F[\mathcal Y]\oplus\operatorname{im}(\operatorname{ad}\tilde A)$.
\item[(e)] The projection $\operatorname{pr}$ onto $F[\mathcal Y]$ along $\operatorname{im}(\operatorname{ad}\tilde A)$ is $F$-linear and preserves the shift; for $X\in\mathfrak{gl}^{(m)}$, $1\leq m\leq5$,
\[
\operatorname{pr}(X)=\tfrac15\lambda_m(X)\,\tilde A^m=\tfrac15\Bigl[\sum_a\xi_a\Bprod^{(5-m)}_{a+m}\Bigr]\mathcal Y^{5-m}.
\]
Equivalently, $\operatorname{pr}(X)$ is the unique element of $F[\mathcal Y]$ with $\tr(\operatorname{pr}(X)Z)=\tr(XZ)$ for all $Z\in F[\mathcal Y]$.
\item[(f)] Let $F=\Bbbk\lser u$ and let $\tilde A$ have the one-pole shape. Then $\mathcal Y\in\mathfrak{gl}_5(\Bbbk\fps u)$; an element $\sum_{j=0}^4w_j\mathcal Y^j$ ($w_j\in\Bbbk\lser u$) lies in $\mathfrak{gl}_5(\Bbbk\fps u)$ if and only if every $w_j$ lies in $\Bbbk\fps u$; the commutant of $\tilde A$ in $\mathfrak{gl}_5(\Bbbk\fps u)$ is the ring $\Bbbk\fps u[\mathcal Y]:=\bigoplus_{j=0}^4\Bbbk\fps u\,\mathcal Y^j$; and for $X\in\mathfrak{gl}_5(\Bbbk\fps u)$ the coefficient of $\mathcal Y^j$ in $\operatorname{pr}(X)$ lies in $u^{-1}\Bbbk\fps u$ for $j=1,\dots,4$ and in $\Bbbk\fps u$ for $j=0$.
\item[(g)] (Not used.) $x^5-\Bprod$ is separable, so $\tilde A$ is regular semisimple; in the situation of (f), $x^5-\Bprod$ is irreducible over $\Bbbk\lser u$, so $F[\mathcal Y]$ is a field.
\end{enumerate}
\end{lem}

\begin{proof}
(a) By induction, $\tilde A^{m+1}=\sum_a\Bprod^{(m)}_a\tilde b_{a+m}E_{a,a+m+1}$ and $\Bprod^{(m)}_a\tilde b_{a+m}=\Bprod^{(m+1)}_a$; for $m=5$, $E_{a,a+5}=E_{aa}$. $\tilde A\sum_c\tilde b_c^{-1}E_{c+1,c}=\sum_aE_{aa}$ gives $\mathcal Y$, and $\mathcal Y^j$ follows in the same way. Reindexing, $\mathcal Y^{5-m}=\sum_a(\Bprod^{(5-m)}_{a+m})^{-1}E_{a,a+m}$, and $\Bprod^{(m)}_a\Bprod^{(5-m)}_{a+m}=\Bprod$, so $\mathcal Y^{5-m}=\Bprod^{-1}\tilde A^m$. In the Leibniz expansion of $\det(x1-\tilde A)$ a permutation $\pi$ contributes only if $\pi(a)\in\{a,a+1\}$ for all $a$; if $\pi(a)=a+1$ for one $a$, then $\pi(a+1)\neq a+1$, so $\pi(a+1)=a+2$, and so on: $\pi$ is the $5$-cycle $a\mapsto a+1$, which is even and contributes $\prod_a(-\tilde b_a)=-\Bprod$; otherwise $\pi$ is the identity.

(b) $\tilde AX=\sum_a\tilde b_a\xi_{a+1}E_{a,a+1+m}$ and $X\tilde A=\sum_a\xi_a\tilde b_{a+m}E_{a,a+m+1}$.

(c) By (b), $\operatorname{ad}\tilde A$ maps $\mathfrak{gl}^{(m)}$ to $\mathfrak{gl}^{(m+1)}$ and the sum of the $\mathfrak{gl}^{(m)}$ is direct, so $X$ commutes with $\tilde A$ if and only if each shift component does. For $X\in\mathfrak{gl}^{(m)}$ this means $\xi_{a+1}=\xi_a\tilde b_{a+m}/\tilde b_a$ for all $a$: a solution is determined by $\xi_0$, and $\xi_a=\Bprod^{(m)}_a$ is a non-zero solution, since $\Bprod^{(m)}_{a+1}\tilde b_a=\Bprod^{(m)}_a\tilde b_{a+m}$. So $\ker(\operatorname{ad}\tilde A)\cap\mathfrak{gl}^{(m)}=F\tilde A^m$ and $\ker(\operatorname{ad}\tilde A)=\bigoplus_{m=1}^5F\tilde A^m=F[\tilde A]$. By (a), $\tilde A^m=\Bprod\mathcal Y^{5-m}$ and $\mathcal Y=\Bprod^{-1}\tilde A^4$, so $F[\tilde A]=F[\mathcal Y]=\bigoplus_{j=0}^4F\mathcal Y^j$.

(d) For $Z\in\ker(\operatorname{ad}\tilde A)$ and any $Y$, $\tr(Z[\tilde A,Y])=\tr(\tilde AZY)-\tr(Z Y\tilde A)=0$. Take $Z=\tilde A^{5-m}$ and $Y\in\mathfrak{gl}^{(m-1)}$; for $X\in\mathfrak{gl}^{(m)}$, (a) gives $X\tilde A^{5-m}=\sum_a\xi_a\Bprod^{(5-m)}_{a+m}E_{aa}$, so $\tr(X\tilde A^{5-m})=\sum_a\xi_a\Bprod/\Bprod^{(m)}_a=\Bprod\,\lambda_m(X)$. Hence $\lambda_m$ vanishes on $\operatorname{ad}\tilde A(\mathfrak{gl}^{(m-1)})$. This image has dimension $5-1=4$ by (c), and $\lambda_m\neq0$ because $\lambda_m(\tilde A^m)=\sum_a\Bprod^{(m)}_a/\Bprod^{(m)}_a=5\neq0$; so the image is $\ker\lambda_m$, and $\mathfrak{gl}^{(m)}=F\tilde A^m\oplus\ker\lambda_m$. By (b), $\operatorname{im}(\operatorname{ad}\tilde A)=\bigoplus_m\operatorname{ad}\tilde A(\mathfrak{gl}^{(m-1)})$, and summing over $m$ gives $\mathfrak{gl}_5(F)=F[\tilde A]\oplus\operatorname{im}(\operatorname{ad}\tilde A)$.

(e) $X=\frac15\lambda_m(X)\tilde A^m+\bigl(X-\frac15\lambda_m(X)\tilde A^m\bigr)$ with the second term in $\ker\lambda_m$; the second form uses $\tilde A^m=\Bprod\mathcal Y^{5-m}$ and $\Bprod/\Bprod^{(m)}_a=\Bprod^{(5-m)}_{a+m}$. $X-\operatorname{pr}(X)$ lies in $\operatorname{im}(\operatorname{ad}\tilde A)$, which is orthogonal to $F[\mathcal Y]$ under $(X,Z)\mapsto\tr(XZ)$ by the first line of (d). The trace form is non-degenerate on $F[\tilde A]$: $\tilde A^n$ has shift $n$, so $\tr\tilde A^n=0$ unless $5\mid n$, $\tr1=5$ and $\tr\tilde A^5=5\Bprod$; in the basis $1,\tilde A,\dots,\tilde A^4$ the Gram matrix has $5$ at $(0,0)$, $5\Bprod$ at $(i,5-i)$ for $i=1,\dots,4$ and $0$ elsewhere, with determinant $5(5\Bprod)^4\neq0$. This gives the uniqueness.

(f) The entries of $\mathcal Y$ are the $\tilde b_a^{-1}$, which are integral ($\tilde b_1^{-1}=u(u\tilde b_1)^{-1}$). By (a), $\mathcal Y^j$ is supported on the positions $(a+j,a)$; for $j=0,\dots,4$ these five sets partition the $25$ positions, so the entry $(a+j,a)$ of $\sum_jw_j\mathcal Y^j$ is $w_j(\Bprod^{(j)}_a)^{-1}$. For $j\leq4$ the window $\{a,\dots,a+j-1\}$ has at most four elements and contains the pole index $1$ at most once, so $\Bprod^{(j)}_a$ is a unit or $u^{-1}$ times a unit and $(\Bprod^{(j)}_a)^{-1}$ is integral; for $a=2$ the window lies in $\{2,3,4,0\}$ and $(\Bprod^{(j)}_2)^{-1}$ is a unit. Hence the element is integral if and only if every $w_j$ is. $\Bbbk\fps u[\mathcal Y]$ is a ring because $\mathcal Y^5=\Bprod^{-1}1$ and $\Bprod^{-1}\in u\Bbbk\fps u$; by (c) it is the integral commutant. For $X\in\mathfrak{gl}^{(m)}\cap\mathfrak{gl}_5(\Bbbk\fps u)$, $1\leq m\leq4$, the coefficient of $\mathcal Y^{5-m}$ in $\operatorname{pr}(X)$ is $\frac15\sum_a\xi_a\Bprod^{(5-m)}_{a+m}$, a sum of products of an integral $\xi_a$ with at most four consecutive weights, at most one of them $\tilde b_1$, so it lies in $u^{-1}\Bbbk\fps u$; for $m=5$ it is $\frac15\tr X$, which is integral.

(g) $\gcd(x^5-\Bprod,5x^4)=1$ since $\Bprod\neq0$, so $\tilde A$ has five distinct eigenvalues over an algebraic closure. In the situation of (f) write $\Bprod=\Bprod^\circ/u$ with $\Bprod^\circ:=u\Bprod\in\Bbbk\fps u^\times$ (a local letter); with $x=1/X$, $x^5-\Bprod=-\Bprod^\circ u^{-1}X^{-5}\bigl(X^5-u(\Bprod^\circ)^{-1}\bigr)$, and $X^5-u(\Bprod^\circ)^{-1}$ is an Eisenstein polynomial at the prime $u$ of $\Bbbk\fps u$.
\end{proof}

% Section 3, part 4 of 7: lifting, the scalar equation, proof of Theorem U, the leading term, the Z/5-degree, graded gauges.

\begin{lem}[Lifting]
\label{lem:gau-lift}
Let $\tilde A$ have the one-pole shape, let $X\in\mathfrak{gl}_5(\Bbbk\fps u)$ have shift $m+1$ with $0\leq m\leq4$, and suppose $\operatorname{pr}(X)=0$. Then there is $\mathcal Z\in\mathfrak{gl}_5(\Bbbk\fps u)$ of shift $m$ with $[\tilde A,\mathcal Z]=X$, and every solution $\mathcal Z'\in\mathfrak{gl}_5(\Bbbk\fps u)$ of $[\tilde A,\mathcal Z']=X$ is $\mathcal Z+w$ with $w\in\Bbbk\fps u[\mathcal Y]$. If $\tilde A$ and $X$ have rational coefficients, $\mathcal Z$ can be chosen rational.
\end{lem}

\begin{proof}
Write $X=\sum_a\xi_aE_{a,a+m+1}$ and look for $\mathcal Z=\sum_a\rho_a\Bprod^{(m)}_aE_{a,a+m}$ with $\rho_a\in\Bbbk\lser u$ (local letters). By \cref{lem:gau-comm}~(b) and $\tilde b_a\Bprod^{(m)}_{a+1}=\Bprod^{(m)}_a\tilde b_{a+m}=\Bprod^{(m+1)}_a$,
\[
[\tilde A,\mathcal Z]=\sum_a(\rho_{a+1}-\rho_a)\Bprod^{(m+1)}_aE_{a,a+m+1},
\]
so the equation is $\rho_{a+1}-\rho_a=\sigma_a:=\xi_a/\Bprod^{(m+1)}_a$ for $a\in\bbZ/5$. It is solvable if and only if $\sum_a\sigma_a=\lambda_{m+1}(X)=0$, i.e.\ if and only if $\operatorname{pr}(X)=0$ (\cref{lem:gau-comm}~(e)). Each $\sigma_a$ is integral, and $\sigma_a(0)=0$ whenever the window $\{a,\dots,a+m\}$ contains the pole index $1$, because then $\Bprod^{(m+1)}_a$ is $u^{-1}$ times a unit; for $m+1=5$ every window contains $1$. Put $\rho_1:=0$ and $\rho_{a+1}:=\rho_a+\sigma_a$ for $a=1,2,3,4$; the relation for $a=0$ then holds because $\sum_a\sigma_a=0$. All $\rho_a$ are integral. The entry $\rho_a\Bprod^{(m)}_a$ is integral if the window $\{a,\dots,a+m-1\}$ avoids $1$, and if it contains $1$ (that is, for the $m$ indices $a=2-m,\dots,1$) it is integral provided $\rho_a(0)=0$. For $a=1$ this holds by the choice $\rho_1=0$; for $a=2-m,\dots,0$ one has $\rho_a=\rho_1-\sum_{c=a}^0\sigma_c$, and every window $\{c,\dots,c+m\}$ with $2-m\leq c\leq0$ contains $1$, so $\rho_a(0)=0$. Hence $\mathcal Z$ is integral. Two integral solutions differ by an element of $\ker(\operatorname{ad}\tilde A)\cap\mathfrak{gl}_5(\Bbbk\fps u)=\Bbbk\fps u[\mathcal Y]$ (\cref{lem:gau-comm}~(c), (f)). The construction uses only field operations on the data.
\end{proof}

\begin{lem}[The scalar equation]
\label{lem:gau-scal}
Let $A$, $\tilde A$ have the one-pole shape and satisfy (EP).
\begin{enumerate}
\item[(a)] For $w=\sum_{j=0}^4w_j\mathcal Y^j$ with $w_j\in\Bbbk\lser u$,
\[
\operatorname{pr}\bigl(w'+\operatorname{dlog}\Theta_U\,w\bigr)=\sum_{j=0}^4\bigl(w_j'+\chi^U_jw_j\bigr)\mathcal Y^j,\qquad \chi^U_j:=-\frac j5\,\frac{\Bprod'}{\Bprod}+\frac15\tr(\operatorname{dlog}\Theta_U).
\]
In particular $\operatorname{pr}((\mathcal Y^j)')=-\frac j5(\Bprod'/\Bprod)\mathcal Y^j$ and $\operatorname{pr}(\operatorname{dlog}\Theta_U\,\mathcal Y^j)=\frac15\tr(\operatorname{dlog}\Theta_U)\,\mathcal Y^j$. Moreover $\chi^U_j=\frac j{5u}+\chi^\circ_j$ with $\chi^\circ_j:=-\frac j5(\Bprod^\circ)'/\Bprod^\circ+\frac15\tr(\operatorname{dlog}\Theta_U)\in\Bbbk\fps u$, where $\Bprod^\circ:=u\Bprod\in\Bbbk\fps u^\times$.
\item[(b)] For $Z\in\mathfrak{gl}_5(\Bbbk\fps u)$ the equation $\operatorname{pr}(w'+\operatorname{dlog}\Theta_U\,w)=-\operatorname{pr}(Z)$ has solutions $w\in\Bbbk\fps u[\mathcal Y]$; their components $w_1,\dots,w_4$ are unique, and $w_0$ is unique up to $w_0\mapsto w_0+\cn\,\mathfrak d$, $\cn\in\Bbbk$.
\item[(c)] For $Z=0$ the solutions in $\Bbbk\fps u[\mathcal Y]$ are exactly the $\cn\,\mathfrak d\,1$, $\cn\in\Bbbk$.
\item[(d)] If the data are rational, so are the solutions for $\cn\in\bbQ$; $\mathfrak d\in1+u\bbQ\fps u$.
\end{enumerate}
\end{lem}

\begin{proof}
(a) $(\mathcal Y^j)'=\sum_a((\Bprod^{(j)}_a)^{-1})'E_{a+j,a}$ has shift $5-j$; written as $\sum_{a'}\xi_{a'}E_{a',a'+5-j}$ it has $a'=a+j$ and $\xi_{a'}=((\Bprod^{(j)}_{a'-j})^{-1})'$. By \cref{lem:gau-comm}~(e) with $m=5-j$ (and $m=5$ for $j=0$),
\[
\operatorname{pr}((\mathcal Y^j)')=\frac15\Bigl[\sum_a\bigl((\Bprod^{(j)}_a)^{-1}\bigr)'\Bprod^{(j)}_a\Bigr]\mathcal Y^j=-\frac15\Bigl[\sum_a\frac{(\Bprod^{(j)}_a)'}{\Bprod^{(j)}_a}\Bigr]\mathcal Y^j .
\]
Here $(\Bprod^{(j)}_a)'/\Bprod^{(j)}_a=\sum_{i=0}^{j-1}\tilde b_{a+i}'/\tilde b_{a+i}$, and as $a$ runs over $\bbZ/5$ each index $c$ occurs for exactly the $j$ pairs $(a,i)=(c-i,i)$; so the sum is $j\sum_c\tilde b_c'/\tilde b_c=j\Bprod'/\Bprod$. Next, $\operatorname{dlog}\Theta_U\,\mathcal Y^j=\sum_a\frac{(d^U_{a+j})'}{d^U_{a+j}}(\Bprod^{(j)}_a)^{-1}E_{a+j,a}$, so now $\xi_{a'}=\frac{(d^U_{a'})'}{d^U_{a'}}(\Bprod^{(j)}_{a'-j})^{-1}$, and \cref{lem:gau-comm}~(e) gives $\frac15\bigl[\sum_{a'}(d^U_{a'})'/d^U_{a'}\bigr]\mathcal Y^j=\frac15\tr(\operatorname{dlog}\Theta_U)\mathcal Y^j$. As $\operatorname{pr}$ is $\Bbbk\lser u$-linear and fixes $\Bbbk\lser u[\mathcal Y]$, $\operatorname{pr}(w'+\operatorname{dlog}\Theta_Uw)=\sum_j\operatorname{pr}\bigl(w_j'\mathcal Y^j+w_j(\mathcal Y^j)'+w_j\operatorname{dlog}\Theta_U\mathcal Y^j\bigr)=\sum_j(w_j'+\chi^U_jw_j)\mathcal Y^j$. Finally $\Bprod=\Bprod^\circ/u$ with $\Bprod^\circ$ a unit by the one-pole shape, so $\Bprod'/\Bprod=(\Bprod^\circ)'/\Bprod^\circ-u^{-1}$, and $\tr(\operatorname{dlog}\Theta_U)=\sum_c(d^U_c)'/d^U_c$ is integral because the $d^U_c$ are units.

(b) By \cref{lem:gau-comm}~(f), $\operatorname{pr}(Z)=\sum_j\phi^Z_j\mathcal Y^j$ with $u\phi^Z_j\in\Bbbk\fps u$ for $j\geq1$ and $\phi^Z_0\in\Bbbk\fps u$ (local letters). By (a) and the linear independence of the $\mathcal Y^j$ the equation splits into $w_j'+\chi^U_jw_j=-\phi^Z_j$, $j=0,\dots,4$. For $j=1,\dots,4$ multiply by $u$: $uw_j'+(\frac j5+u\chi^\circ_j)w_j=-u\phi^Z_j$. For $w_j=\sum_nw_{j,n}u^n$ the coefficient of $u^n$ reads
\[
\Bigl(n+\frac j5\Bigr)w_{j,n}=-\coef{u^n}(u\phi^Z_j)-\sum_{l=0}^{n-1}\coef{u^{n-1-l}}\chi^\circ_j\cdot w_{j,l},
\]
and $n+j/5$ is a non-zero rational number for every $n\geq0$: so there is exactly one solution $w_j\in\Bbbk\fps u$. For $j=0$, $\chi^U_0=\frac15\tr(\operatorname{dlog}\Theta_U)=-\mathfrak d'/\mathfrak d$, so the equation is $\mathfrak d\,(w_0/\mathfrak d)'=-\phi^Z_0$; its solutions in $\Bbbk\fps u$ are $w_0=\mathfrak d\bigl(\cn-\int_0^u\phi^Z_0/\mathfrak d\,\dd u\bigr)$, $\cn\in\Bbbk$, where $\int_0^u$ is the formal antiderivative without constant term, and the difference $e$ of two solutions satisfies $(e/\mathfrak d)'=0$, so $e\in\Bbbk\mathfrak d$.

(c) is (b) with $\phi^Z=0$: the recursion gives $w_j=0$ for $j\geq1$, and $w_0=\cn\mathfrak d$. (d) The pivots $n+j/5$ are rational, $\det\Theta_U/\det\Theta_U(0)$ lies in $1+u\bbQ\fps u$ and its binomial series $\mathfrak d$ has rational coefficients.
\end{proof}

\begin{proof}[Proof of \cref{thm:gau-U}~(a)]
\emph{Reduction.} Write $P=\Theta_U\check P$; $\check P$ (a local letter) has coefficients in $\mathfrak{gl}_5(\Bbbk\fps u)$ if and only if $P$ has, since $\Theta_U$ is invertible over $\Bbbk\fps u$. As $A\Theta_U=\Theta_U\tilde A$, we have $AP-P\tilde A=\Theta_U(\tilde A\check P-\check P\tilde A)$ and $zP'=z\Theta_U'\check P+z\Theta_U\check P'$, so the gauge equation is equivalent to $z\check P'=[\tilde A,\check P]-z\operatorname{dlog}\Theta_U\,\check P$, i.e.\ to
\begin{equation}
\label{eq:gau-reduced}
[\tilde A,\check P_0]=0,\qquad [\tilde A,\check P_{k+1}]=\check P_k'+\operatorname{dlog}\Theta_U\,\check P_k\quad(k\geq0).
\end{equation}
\emph{Construction.} We construct $\check P_k=\mathcal Z_k+w^{(k)}$ with $\mathcal Z_k\in\mathfrak{gl}_5(\Bbbk\fps u)$ and $w^{(k)}\in\Bbbk\fps u[\mathcal Y]$ such that, for every $k\geq0$, (I$_k$) $[\tilde A,\check P_k]=\check P_{k-1}'+\operatorname{dlog}\Theta_U\check P_{k-1}$ (with $\check P_{-1}:=0$) and (S$_k$) $\operatorname{pr}(\check P_k'+\operatorname{dlog}\Theta_U\check P_k)=0$. For $k=0$ put $\mathcal Z_0:=0$ and $w^{(0)}:=\mathfrak d\,1$: (I$_0$) holds because $w^{(0)}$ commutes with $\tilde A$, and (S$_0$) is \cref{lem:gau-scal}~(c). Given $\check P_k$, the matrix $X_{k+1}:=\check P_k'+\operatorname{dlog}\Theta_U\check P_k$ is integral with $\operatorname{pr}(X_{k+1})=0$ by (S$_k$); since $\operatorname{pr}$ preserves the shift, \cref{lem:gau-lift} applied to each shift component gives an integral $\mathcal Z_{k+1}$ with $[\tilde A,\mathcal Z_{k+1}]=X_{k+1}$. Choose $w^{(k+1)}\in\Bbbk\fps u[\mathcal Y]$ by \cref{lem:gau-scal}~(b) with $Z:=\mathcal Z_{k+1}'+\operatorname{dlog}\Theta_U\mathcal Z_{k+1}$ and any value of the free constant. Then (I$_{k+1}$) holds since $[\tilde A,w^{(k+1)}]=0$, and (S$_{k+1}$) holds since $\operatorname{pr}(\check P_{k+1}'+\operatorname{dlog}\Theta_U\check P_{k+1})=\operatorname{pr}(Z)+\operatorname{pr}\bigl((w^{(k+1)})'+\operatorname{dlog}\Theta_Uw^{(k+1)}\bigr)=0$. By \eqref{eq:gau-reduced}, $P^\natural:=\Theta_U\sum_kz^k\check P_k$ lies in $\Gau(A,\tilde A)$, and $P^\natural_0=\mathfrak d\,\Theta_U$. The only choices are the free constant of \cref{lem:gau-scal}~(b) at each order and the normalisation $\rho_1=0$ inside \cref{lem:gau-lift}, which the next choice of $w^{(k+1)}$ absorbs.
\end{proof}

\begin{lem}[The leading term of a solution of the gauge equation]
\label{lem:gau-lead}
Let $A$, $\tilde A$ have the one-pole shape and satisfy (EP), let $P\in\Gau(A,\tilde A)$ and $n\geq0$ with $P_0=\dots=P_{n-1}=0$. Then $P_n=\cn\,\mathfrak d\,\Theta_U$ for a unique $\cn\in\Bbbk$, namely $\cn=(P_n)_{00}(0)$.
\end{lem}

\begin{proof}
Put $\check P:=\Theta_U^{-1}P$, so $\check P_l=0$ for $l<n$. The coefficient of $z^n$ in the reduced equation \eqref{eq:gau-reduced} is $[\tilde A,\check P_n]=0$ (for $n=0$ this is the first equation), so $\check P_n\in\ker(\operatorname{ad}\tilde A)\cap\mathfrak{gl}_5(\Bbbk\fps u)=\Bbbk\fps u[\mathcal Y]$ (\cref{lem:gau-comm}~(c), (f)). The coefficient of $z^{n+1}$ says that $\check P_n'+\operatorname{dlog}\Theta_U\check P_n=[\tilde A,\check P_{n+1}]$ lies in $\operatorname{im}(\operatorname{ad}\tilde A)$, so $\operatorname{pr}(\check P_n'+\operatorname{dlog}\Theta_U\check P_n)=0$. By \cref{lem:gau-scal}~(c), $\check P_n=\cn\,\mathfrak d\,1$. As $d^U_0=1$ and $\mathfrak d(0)=1$, $\cn=(P_n)_{00}(0)$.
\end{proof}

\begin{proof}[Proof of \cref{thm:gau-U}~(b) and (d)]
(b) Let $P^\natural\in\Gau(A,\tilde A)$ with $P^\natural_0=\mathfrak d\,\Theta_U$; the gauge of (a) is one. \emph{Injectivity:} if $\cn\neq0$ has order $n$ in $z$, then $\coef{z^n}(P^\natural\cn)=\cn_nP^\natural_0=\cn_n\mathfrak d\Theta_U\neq0$. \emph{Surjectivity:} let $P\in\Gau(A,\tilde A)$. We define $\cn_0,\cn_1,\dots\in\Bbbk$ recursively. Suppose $R^{(n)}:=P-P^\natural(\cn_0+\cn_1z+\dots+\cn_{n-1}z^{n-1})$ is $0$ modulo $z^n$ (for $n=0$ this is empty); $R^{(n)}\in\Gau(A,\tilde A)$ because $\Gau(A,\tilde A)$ is a $\Bbbk\fps z$-module. By \cref{lem:gau-lead}, $\coef{z^n}R^{(n)}=\cn_n\mathfrak d\Theta_U$ for a unique $\cn_n\in\Bbbk$. Then $R^{(n+1)}=R^{(n)}-\cn_nz^nP^\natural$ lies in $\Gau(A,\tilde A)$ and $\coef{z^n}R^{(n+1)}=\cn_n\mathfrak d\Theta_U-\cn_nP^\natural_0=0$, so $R^{(n+1)}\equiv0$ modulo $z^{n+1}$. Put $\cn:=\sum_n\cn_nz^n$. For every $l$ and every $n>l$, $\coef{z^l}(P-P^\natural\cn)=\coef{z^l}R^{(n)}-\coef{z^l}P^\natural(\cn-\sum_{m<n}\cn_mz^m)=0$. Hence $P=P^\natural\cn$. Its $z^0$-term is $\cn(0)\mathfrak d\Theta_U$, invertible if and only if $\cn(0)\neq0$. If $P=P^\natural\cn_P$ is a regular gauge, every regular gauge $P^\natural\cn'$ equals $P\cdot(\cn_P^{-1}\cn')$ with $\cn_P^{-1}\cn'\in\Bbbk\fps z^\times$, and conversely. For the pair $(\tilde A,\tilde A)$ one has $\Theta_U=1$ and $\mathfrak d=1$, and $1\in\Gau(\tilde A,\tilde A)$; so $\Gau(\tilde A,\tilde A)=\Bbbk\fps z\cdot1$.

(d) Every object used in the proof of (a) is obtained from the $b_a$, $\tilde b_a$ by field operations in $\bbQ\lser u$, the derivative and the integration $\int_0^u$ of power series: $\Theta_U$, $\operatorname{dlog}\Theta_U$, $\mathcal Y$, $\Bprod^{(m)}_a$, $\Bprod^\circ$, $\operatorname{pr}$ (\cref{lem:gau-comm}~(e)), $\mathfrak d$, the lift of \cref{lem:gau-lift} and the solutions of \cref{lem:gau-scal}~(b) (\cref{lem:gau-scal}~(d)). With every free constant chosen in $\bbQ$, the construction therefore runs inside $\mathfrak{gl}_5(\bbQ\fps u)\fps z$. For a field $\Bbbk\supseteq\bbQ$ the proofs of \cref{lem:gau-comm,lem:gau-lift,lem:gau-scal,lem:gau-lead} hold verbatim over $\Bbbk$, so the solutions over $\Bbbk$ are $P^\natural\Bbbk\fps z$ by (b). If $P=P^\natural\cn$ has rational coefficients, the recursion of the surjectivity proof computes $\cn_n=(\coef{z^n}R^{(n)})_{00}(0)$ from rational data, so $\cn\in\bbQ\fps z$.
\end{proof}

\begin{lem}[The $\bbZ/5$-degree of the gauge equation]
\label{lem:gau-degree}
For $P=\sum_kz^kP_k$ with $P_k\in\mathfrak{gl}_5(\Bbbk\lser u)$ and $\delta\in\bbZ/5$ put $P^{[\delta]}:=\sum_kz^k\sum_{(a,b):\,k+b-a\equiv\delta}(P_k)_{ab}E_{ab}$, so that $P=\sum_\delta P^{[\delta]}$, and $P$ is graded if and only if $P=P^{[0]}$. Let $\operatorname{ge}(P):=zP'-AP+P\tilde A$. Then $\operatorname{ge}(P^{[\delta]})=\operatorname{ge}(P)^{[\delta+1]}$; in particular every degree part of a solution of the gauge equation is a solution.
\end{lem}

\begin{proof}
Give $z^kf(u)E_{ab}$ the degree $k+b-a$ modulo $5$. The derivative in $u$ does not change it and the factor $z$ adds $1$, so $zP'$ raises the degree of each homogeneous term by $1$. $A\,(fz^kE_{cd})=b_{c-1}fz^kE_{c-1,d}$ has degree $k+d-c+1$, and $(fz^kE_{cd})\tilde A=\tilde b_dfz^kE_{c,d+1}$ has degree $k+d+1-c$. So $\operatorname{ge}$ maps homogeneous series of degree $\delta$ to homogeneous series of degree $\delta+1$, and by additivity $\operatorname{ge}(P)^{[\delta+1]}=\operatorname{ge}(P^{[\delta]})$.
\end{proof}

\begin{prop}[Graded gauges]
\label{prop:gau-graded}
Let $A$, $\tilde A$ have the one-pole shape and satisfy (EP).
\begin{enumerate}
\item Graded regular gauges exist; they can be chosen rational if all $b_a$, $\tilde b_a$ are rational.
\item If $P$ is a graded regular gauge, the graded regular gauges are exactly the $P\cn(z)$ with $\cn\in\Bbbk\fps{z^5}^\times$.
\item For a graded regular gauge $P$, $P_{11}(0,z)=\sum_{m\geq0}(P_{5m})_{11}(0)\,z^{5m}$ is a unit of $\Bbbk\fps{z^5}$.
\item If $\zeta\in\Bbbk$, then $\operatorname{Rot}_iP=\sum_\delta\zeta^{-i\delta}P^{[\delta]}$; so $P$ is graded if and only if $\operatorname{Rot}_1P=P$, if and only if $\operatorname{Rot}_iP=P$ for all $i$, and $\frac15\sum_i\operatorname{Rot}_iP=P^{[0]}$. Moreover $A\,\zeta^{i(\mathrm{gr}+1)}=\zeta^i\zeta^{i(\mathrm{gr}+1)}A$, and the same for $\tilde A$; $\operatorname{Rot}_i$ maps solutions of $z\de_uY=AY$ to solutions, and $\operatorname{Rot}_i(PY)=(\operatorname{Rot}_iP)(\operatorname{Rot}_iY)$, $\operatorname{Rot}_i(Y\mathsf d(z))=\operatorname{Rot}_i(Y)\,\mathsf d(\zeta^{-i}z)$ for scalar series $\mathsf d$.
\end{enumerate}
\end{prop}

\begin{proof}
(i) Let $P^\natural$ be as in \cref{thm:gau-U}~(a), rational in the rational case by (d). Its $z^0$-term $\mathfrak d\Theta_U$ is diagonal, of degree $0$. By \cref{lem:gau-degree}, $P^{\mathrm{gr}}:=(P^\natural)^{[0]}$ solves the gauge equation, and $P^{\mathrm{gr}}_0=P^\natural_0$ is invertible. (ii) By \cref{thm:gau-U}~(b) the regular gauges are the $P\cn(z)$, $\cn\in\Bbbk\fps z^\times$. Write $\cn=\sum_{r\in\bbZ/5}\cn^{(r)}$ with $\cn^{(r)}:=\sum_{n\equiv r}\cn_nz^n$. As $P$ has degree $0$ and $z^n$ adds $n$ to the degree, $(P\cn)^{[r]}=P\cn^{(r)}$. So $P\cn$ is graded if and only if $P\cn^{(r)}=0$ for $r\neq0$, if and only if $\cn^{(r)}=0$ for $r\neq0$, because $\cn\mapsto P\cn$ is injective ($P=P^\natural\cn_P$ with $\cn_P\neq0$, and $\Bbbk\fps z$ is a domain); that is, if and only if $\cn\in\Bbbk\fps{z^5}$. (iii) For graded $P$ the entry $(P_k)_{11}$ vanishes unless $1\equiv1+k$, so $P_{11}(0,z)\in\Bbbk\fps{z^5}$, and its constant term is $(P_0)_{11}(0)=\cn(0)\mathfrak d(0)d^U_1(0)\neq0$ by \cref{thm:gau-U}~(b). (iv) The entry $(a,b)$ of $\coef{z^k}\operatorname{Rot}_iP$ is $\zeta^{i(a+1)-ik-i(b+1)}(P_k)_{ab}=\zeta^{-i(k+b-a)}(P_k)_{ab}$, which gives the first assertion; the equivalences follow because the decomposition $P=\sum_\delta P^{[\delta]}$ is direct, $\zeta^{-i\delta}\neq1$ for $i,\delta\not\equiv0$, and $\sum_i\zeta^{-i\delta}=5\,\delta_{\delta,0}$. The entry $(a,a+1)$ of $\zeta^{i(\mathrm{gr}+1)}A\zeta^{-i(\mathrm{gr}+1)}$ is $\zeta^{i(a+1)}b_a\zeta^{-i(a+2)}=\zeta^{-i}b_a$, which is the commutation rule. Hence, for a solution $Y$, $z\de_u(\operatorname{Rot}_iY)=\zeta^i\zeta^{i(\mathrm{gr}+1)}\bigl[(\zeta^{-i}z)\de_uY\bigr](u,\zeta^{-i}z)=\zeta^i\zeta^{i(\mathrm{gr}+1)}A\,Y(u,\zeta^{-i}z)=A\operatorname{Rot}_iY$. The two product rules are immediate from the definitions.
\end{proof}

\begin{rmk}[What the proof of Theorem~U uses; the role of the pole]
\label{rmk:gau-Uremarks}
The proof of \cref{thm:gau-U} uses only the one-pole shape, (EP) and the characteristic $0$ (in $\lambda_m(\tilde A^m)=5\neq0$, in the pivots $n+j/5\neq0$, and for $\mathfrak d$ and the constants of $\dd/\dd u$). The pole of $\tilde b_1$ enters exactly through the integrality statements, \cref{lem:gau-comm}~(f) and the window argument of \cref{lem:gau-lift}, and through the term $j/(5u)$ of $\chi^U_j$. For two systems without a pole (all ten weights units, with (EP)) the integrality statements hold trivially and $\chi^U_j\in\Bbbk\fps u$ for every $j$, so in the proof of \cref{lem:gau-scal}~(b) every equation $w_j'+\chi^U_jw_j=-\phi^Z_j$, $j=0,\dots,4$, has the solutions $\mathcal E_j\bigl(\mathsf a_j-\int_0^u\phi^Z_j/\mathcal E_j\,\dd u\bigr)$, $\mathsf a_j\in\Bbbk$, with $\mathcal E_j:=\exp(-\int_0^u\chi^U_j\,\dd u)$ (local letters). Then the proof of \cref{lem:gau-lead} gives the leading terms $\Theta_U\sum_j\mathsf a_j\mathcal E_j\mathcal Y^j$, the construction of \cref{thm:gau-U}~(a) started from $w^{(0)}:=\mathcal E_j\mathcal Y^j$ gives a solution $P^{(j)}$ for each $j$, and the proof of \cref{thm:gau-U}~(b), with the five leading terms in place of one, shows that $P^{(0)},\dots,P^{(4)}$ form a basis of $\Gau(A,\tilde A)$ over $\Bbbk\fps z$. So the solution module has rank $5$, and $P_{11}(0,z)=1$ no longer fixes a gauge (see \cref{sec:evi-gauge} for a computed example). This is not used.
\end{rmk}

\input{sections/s03_gauge_5}



